\documentclass[openany,nobib]{tufte-book}

\hypersetup{colorlinks}

\newcommand{\subtitle}[1]{\gdef\thesubtitle{#1}}

\renewcommand{\maketitlepage}{%
{%
\begin{fullwidth}%
\fontsize{18}{20}\selectfont\par\noindent
  \textcolor{darkgray}{\it\thanklessauthor}%
\vspace{11.5pc}%
\fontsize{48}{40}\selectfont\par\noindent
  \textcolor{darkgray}{\thanklesstitle}%
\vspace{0.5pc}%
{\fontsize{26}{40}\selectfont\par\noindent
  \textcolor{darkgray}{\thesubtitle}}%
\vfill%
\fontsize{14}{16}\selectfont\par\noindent{\it\thanklesspublisher}%
\end{fullwidth}%
}
\thispagestyle{empty}%
\clearpage%
}

\title{Holon Programming}
\subtitle{A Survey}
\author{Pierre-Arnoul de Marneffe}
\publisher{December 1973}

\usepackage[greek,english]{babel}
% no Palatino Greek in LGR: substitute CM Greek explicitly
\DeclareFontFamily{LGR}{pplj}{}
\DeclareFontShape{LGR}{pplj}{m}{n}{<->ssub*cmr/m/n}{}
\usepackage{csquotes}

\usepackage{tikz}
\usetikzlibrary{positioning,shapes.geometric,matrix,angles,arrows.meta}
\tikzstyle{every picture}+=[font=\sffamily]
\usepackage[edges]{forest}

% program keywords; indent in units of 2mm
\newcommand{\kw}[1]{{\bf #1}}
\newcommand{\ind}[1]{\-\hspace{\dimexpr 2mm*#1\relax}}
\newcommand{\spchar}{'\textvisiblespace\,'}
\newcommand{\na}{\quad--}% used in tables for N/A cells

% BNF
\newcommand{\sym}[1]{$\langle$#1$\rangle$}
\def\is{$::=$\ }
\def\alt{$|$\ }

\usepackage{tcolorbox}
\tcbuselibrary{breakable}
\tcbset{width=1.1\textwidth, halign=justify, colframe=lightgray,
  boxrule=0.4pt, center, breakable, colback=white}

\usepackage[style=verbose,citepages=omit,dashed=false,maxcitenames=1]{biblatex}
\ExecuteBibliographyOptions{citetracker=true}
\addbibresource{holon.bib}
\setlength\bibitemsep{1.2\itemsep}

\begin{document}

% Front matter
\frontmatter

% r.1 blank page
%\blankpage


% r.3 full title page
\maketitle

%%% Forward
\chapter*{Foreword}

In this survey of the Holon Programming Concept,
we shall try to collect together all the different
notes we have written about this problem. All the
texts about Holon Programming were written in French\autocite{marneffe1972}
(except the text for the {\sc ics} presentation)\autocite{marneffe1973}, and
that precludes a diffusion of the concept, and at the
same time, a welcome feedback of criticisms and
comments to the author.

This survey was written during a stay at the IBM Federal
Systems Center (Gaithersburg, Maryland). This stay was
supported by the Belgian Fund for Scientific Research
(F.N.R.S) with the help of the IBM-Belgium Donation.
\clearpage

%*******************************************************
% Abstract
%*******************************************************
%\renewcommand{\abstractname}{Abstract}
\pdfbookmark[1]{Abstract}{Abstract}
\begingroup
\let\clearpage\relax
\let\cleardoublepage\relax
\let\cleardoublepage\relax

\chapter*{Abstract}
The ``Holon Programming Language'' is designed to force
the programmer to design programs according to a
top-down methodology (structured programming). But, the
language is designed to generate well-structured and
efficient programs. The holon concept is explained; the
Holon Language Syntax is described; and some further
developments are presented: stretching the concept to
concurrent processes, and the development of a holon
operating system for a virtual memory computer.

\endgroup

\vfill


% r.5 contents
\tableofcontents

\listoffigures

%%
% Start the main matter (normal chapters)
\mainmatter

\chapter{Letter to the Author from Knuth}
\label{ch:knuth-letter}

\hfill April 1, 1974

\noindent Prof. Pierre-Arnoul de Marneffe\\
\noindent Universit\'e de Li\`ege\\
\noindent Service d'Informatique\\
\noindent Avenue des Tilleuls 59\\
\noindent {\sc b}-4000 Li\`ege, Belgium\\

\bigskip

\noindent Dear Prof. de Marneffe:

\medskip

Thank you very much for sending me your survey of Holon Programming. I
especially enjoyed your references to the non-computer literature (Koestler,
Bernard-Shaw, Shanley, Mount Vernon, etc.) since computer scientists need to
avoid insularity.

I believe you are making important strides toward the development of a new
programming language. There still remain some unclear areas but you are
obviously addressing the correct issues; the next thing to do (it seems to me)
is to program  several hundred examples!

For related reading I would suggest that you carefully study Ole-Johan Dahl's
papers on SIMULA since his class concept is so close to the holon concept. Also
I have just heard that Brian Randell of Newcastle has been working on a
so-called PEARL system.

I found your report could have been improved if you had worked entirely with
tree structures instead of converting to binary trees. The original tree
structure is what is really relevant, and the Dewey notation for such structure
is more directly suited to the operations you discuss. The binary tree is only
a machine-oriented representation of the basic concept, the discussion should
stay at a higher level.

Secondly, I found the report too preoccupied with details of implementation.
The people by whom it is most important that this report be read are either
able to visualize easily how to implement this sort of system, or else they are
people who are not likely to care how it's implemented as long as it's handled
sensibly. The important thing to stress is rather the conceptual issues of how
holon programming differs from and improves on today's languages.

The example didn't come until page 100,\marginnote{page~\pageref{ch:prog} of
this document} while I expect most readers would have preferred to see it
immediately. Since the program is almost self-explanatory, you can let it
explain the language at the same time (integration of functions!).

Ideally there should be more examples of course.

The one example raises some interesting issues since the individual holons
don't quite state their assumptions. In the very first holon, for example, it
is not at all clear why you `find first word starting character' instead of
going right into `find a word \kw{etc}.' You must already have made a decision
(a) that you wouldn't assume the text begins with a nonblank, (b) that there is
going to be at this level an element of data representing the last-read
character, (c) that the `find a word' routine will already have its first
character in hand, and (d) that there is no need to test for a message that has
no words (only a full stop). As I recall when I was solving that problem, it
took me a good five minutes to reach these decisions, during which I must have
considered lots of alternatives. Once this step was made the rest of the
program flowed naturally. My questions are: Where should we state these
assumptions? Shouldn't we mention the existence of data representing the
last-read character, even though we don't want to specify its detailed
structure until later?

These issues seem to arise repeatedly and I haven't a first conclusion about
what we ought to do. That's why I suggest working out hundreds of examples, as
being the best kind of eating to prove the pudding at this stage. On the other
hand creating the holon implementation itself is equivalent to working out
quite a few examples.

Thanks again for showing me your stimulating work. I myself must get on with
the writing of volume 4 of my series, so I have little energy to devote to the
development of languages, but I will do my best to see that other people
working in the area are kept informed of what you are doing.

\medskip

\hfill Sincerely,

\bigskip

\hfill Donald E. Knuth

\hfill Professor

\bigskip
\bigskip

\noindent P.S. Is there a place in Belgium whose postal code is {\sc b}-6700
like the Burroughs computer?


% v.2 epigraphs
\newpage\thispagestyle{empty}


\vspace*{\fill}

\begin{flushright}{
    Your Royal Academician thinks he can get the
    style of Giotto without Giotto's beliefs, and
    correct his perspectives in the bargain. Your
    man of letters thinks he can get Bunyan's
    conviction or Shakespear's apprehension,
    especially if he takes care not to split his
    infinitives. And so with your Doctor of Music,
    who, with their collections of discords duly
    prepared and resolved or retarded or anticipated
    in the manner of the great composers think they
    can learn the art of Palestrina from
    Cherubini's treatise.} \\ \medskip
    --- {\sc \large george bernard shaw} \\
    {\sl Letter to Arthur Bingham Walkley}
\end{flushright}

%************************************************
\chapter[Structured Programming and Present Programming Languages]{Structured
Programming and\\Present Programming Languages}\label{ch:introduction}
%************************************************

It was contended by Dijkstra that: ``\dots the art of programming is the art of
organizing complexity, of mastering multitude and avoiding its bastard chaos as
effectively as possible\dots.''\autocite{dijkstra1971short} His letter to the
{\sc acm}'s Editor about the ``harmful GOTO
statement''\autocite{dijkstra1968letters} was already an attempt to avoid
``chaos'' in programming. A full explanation of a methodology to design
well-formed programs appeared for the first time in his ``Notes on Structured
Programming.''\autocite{Dijkstra1970-ja} It is assumed in the next sections
that the reader is familiar with the notions of structured programming and
top-down programming. (If not, useful references are Dijkstra's above-mentioned
``Notes,'' Wirth,\autocite{wirth1971program}
Woodger,\autocite{woodger1971semantic} and
Mills\autocite{mills1972mathematical}).

Many people use the structured programming methodology on purpose to write
large programs. The unconvinced reader of Dijkstra's ``Notes'' was surely
struck by the publication of practical
results.\autocite{baker1972chief}\autocite{baker1972system} However, the use of
the structured programming methodology with the present programming languages
has, in our opinion, several drawbacks.

A) The programmer cannot use the same language for the design and the
implementation of a program. The design phase of the program is performed by
using a design language which allows only the writing of structured programming
control statements ({\sl ifthenelse}, {\sl dowhile}, {\sl repeatuntil}, {\sl
case}, {\sl sequence}). Those control statements link together a series of
abstract levels, the definition of which is postponed. But, when all the levels
are defined in this way, the programmer must undertake the mapping of the
designed program into an executable one. This last one must be written in a
``machine-understandable'' language. But this mapping is a nontrivial and
error-prone operation.

B) Another problem is the correspondence between each level of the structured
program and an implementation unit for this level. Woodger was the first one,
in my opinion, to point out the lack of a suitable tool for implementing
semantic levels (Woodger's semantic level is similar to Dijkstra's
pearl\footnote{See Section 14 of Dijkstra's {\it Notes}.}).
%\autocite{snowdon}

It is impossible to give to the programmer a ``rule of thumb'' for the relation
between semantic level and implementation unit: suppose you teach a programmer
to structure a program by a stepwise process, introducing a lot of levels---and
that is the right way---you can't advise the programmer to implement each level
as a procedure, because that will generate a very inefficient
program.\autocite{horning1973process} On the other hand, you can't allow the
programmer to follow the rule: create levels as you introduce procedures in
conventional programming. Application of this rule will not produce structured
programs! Because:

1) Structuring programs by procedures leads to programs structured by ``large
steps.'' This way may easily produce incorrect
programs.\autocite{henderson1972experiment}

2) The programmer is influenced by the language which he usually
uses.\autocite{dijkstra1972humble} But, it must be well-understood that he is
influenced not only by the language syntax but also by the idiosyncrasies of
the compiler and the computer he uses. Programmers avoid language statements
which lead to unclear compiler error messages. In his use of procedures, the
programmer is influenced by the relative cost of the procedure overhead on his
computer.

C) When a structured program is translated into a program written in a
conventional language (i.e., not designed for structured programs only), errors
may cut right across the layers, because these are only abstractions (Van der
Poel,\autocite{naur1969software} page 51).

With all these drawbacks in mind, it may seem sensible to design a language
which would force the programmer to design Well-structured programs. The
language must help the programmer to write programs without any computer
architecture in mind. The language structure must permit automatic translation
of a structured program into an efficient one. We shall try to design such a
language by the use of the ``Holon'' concept.

%*****************************************
\chapter{The ``Holon'' Concept}
%*****************************************

We shall show that the close connection of the ``Holon'' concept to the
structured programming principles justifies its use in the language design. A
systematic use of this concept will help to suppress that shortcomings of the
use of the present programming languages for writing structured programs.

The ``Holon'' concept has been introduced in biological and behaviour sciences
by Koestler.\autocite{koestler1968ghost} This concept proceeds from work of
Simon.\autocite{simon1962architecture} It is used for instance to analyse
complex living organisms or complex social systems.

This neologism is from Greek ``holos'' (\textgreek{ὅλος}), i.e., whole, and the
suffix ``-on'' meaning ``part.'' A holon is a ``part of a whole.'' The reader
is forewarned to not mix up the holon concept with the ``module'' one. The
difference will clearly appear, we hope, in the next paragraphs.

According to the holon concept, the complex systems can survive only if they
are holons. A complex system must be a hierarchic ordering of functions
(holons) with the following features:

A) ``Hierarchy'': Each holon is composed by other holons which are
``refinements'' of the former holon. These holons are submitted to some rigid
rules; they perform the ``detail'' operations which, put together, compose the
function of the former holon.

B) ``Tendency to Autonomy'': Each holon has the ability to check its
``environment'' before performing its function; it may modify its behaviour
according to test results. There exists a clear distinction between a ``state
description'' and a ``process description.''\footnote{Ibid.} The holon function
corresponds to the process description; but it checks if the environment upon
which it acts, belongs to its expected state description.

C) ``Tendency to Integration'': The holon integrates with other holons in the
hierarchy according to a flexible strategy. This integration must be understood
as a will to close cooperation with the other holons for the emergence of a
``tougher'' and more efficient component.

A full description of the holon concept appears in Koestler's {\it The Ghost in
the Machine}, chapters three to eight. Only the main features were described in
this section. In the next section, we will display the relation between this
concept and the design of structured programs.

%************************************************
\chapter{Holons and Structured Programming}
%************************************************

Our statement is:

1) The holon model is an attractive one for a programmer using structured
programming;

2) The model may be mapped into a useful and implementable language.

\medskip

\noindent We shall review each of the three holon features to uncover the
corresponding programming features.

\section{Hierarchy}

Anyone, familiar with the structured programming concepts, will immediately
understand the connection between the hierarchic embedding of holons and the
structured programming levels.

\section{Tendency to autonomy}

First of all, this feature must be related to the drawback pointed out by Van
der Poel (see section 1.(c)): errors may cut right across the
layers.\autocite{naur1969software} The ``tendency to autonomy'' may be
understood as checking properties at an appropriate level. These properties
must hold for any correct execution of the program.

Structured Programming may be described in terms of set
theory;\autocite{mills1972mathematical} the correctness of the program is based
upon assertions about functions, the values of which are computed at each
level. For instance, if the programmer is using Hoare's
theorems,\autocite{dijkstra1971short} the correctness is guaranteed if the data
always fit the hypothesis (assertions before each function).

Where? When? How do we do this check? Answers to these three questions will be
given in the following section.

\section{Tendency to integration}

This third holon feature is the most important one. With the help of an
adequate implementation, the programmer will not have to consider anymore the
choice of an implementation unit (open or closed procedure) for each level.
This choice will be automatic and tailored for the computer which will execute
the program. The choice is left to an automatic system because the relative
cost of an open or closed procedure is a machine-dependent parameter. The cost
is expressed as a tradeoff between time and space.

In this way, the programmer designs his program in terms of a set of embedded
functions, and the mapping of the program description in a ``machine'' program
is outside his control. The only thing he knows is that the generated
machine-program will be as efficient as possible for the used computer.

%************************************************
\chapter{Datatest and Invariant}
%************************************************

Now, we return to the three questions about the level properties checking:
Where? When? How?

I must confess that it takes a long time to fully understand the implications
of the questions and the corresponding answers. The holon concept assumes that
a complex system must check its environment to stay alive. It can be made by
the check of the similarity between the real environment and a ``standard''
one.

\section{Where?}

First, consider a holon A. The function of this holon A is made up by the
execution of other holons, for instance: B, C, D. These three holons (B, C, D)
are the first level of the hierarchic structure of the holon A. They may be
composed of other holons, according to the hierarchic embedding of holons.

We say that the execution of holon A is ``sensible'' only if his ``input
environment'' is a ``legal'' one. There exists a set of predicates, the values
of which are {\sl true} if the environment is a ``legal'' one. We call this set
of predicates: the datatest.

Now, consider the three holons (B, C, D), which are the first level of the
hierarchic decomposition of A. Suppose that the execution of A implies the
executions of B, C and D, in this order. After each execution of one of these
holons, the environment of A is modified. Nevertheless, there exists a set of
predicates (upon this environment), the values of which must be {\sl true}
after each execution. If not, the modified environment will not belong to the
range of holon A, after completion of the execution of A. We call this set of
predicates: the invariant.

\section{When?}

When do we check the predicates of the datatest and invariant?

The definitions of these two sets give well-defined positions in the holon: the
datatest at the beginning, and the invariant after each inner holon.

But each set (invariant or datatest) can be segregated in three subsets:

1) A subset of predicates which can be checked at compile-time (for
instance---data types). The extent of this subset depends on the designed
syntax of the language.

2) There are predicates, the truth of which may be proved by the programmer.
The proof can be a formal or informal one relying upon an analysis of the
structure of the program. This means that the programmer can prove the truth of
these predicates, whatever the input to the program.

3) There remains a set of predicates, the values of which will be computed at
run-time.

The main question is: why does the third set of predicates exist? It exists
because the programmer wants to increase the confidence in the program
correctness, the value of the output, and the program robustness.

\section{How?}

The datatest and the invariant will be made up by predicates belonging to the
third set. They are used to trap illegal inputs which will cause havoc. We fear
these inputs because we are unable to trace the consequences of such inputs on
our program, or because we can't decide if this kind of input will or will not
appear during the program lifetime.

For instance, the description of many solutions of some problems are preceded
by the foreword: ``this solution works perfectly provided\dots''; there is room
for a datatest or an invariant.

It may be argued that we don't need a special syntax in the language to
implement this kind of ``fire-wall.'' That is true; but, in our opinion,
casting the program hypothesis in a dedicated syntax unit will catch the
attention of the program reader (and user) to the exact kind of hypothesis
which was made during the program design.

Many readers of previous reports about ``Holon
Programming''\autocite{marneffe1972} have suggested that there is no need for a
distinction between datatest and invariant. I introduced the invariant because
a holon can have a way to verify its output. On the other hand, we can use the
invariant to build a ``fire-wall'' when we merge several predefined holons to
construct a new one. From time to time, we make some kind of bottom-up
construction, even in a top-down design.

The datatest may help by introducing levels in the tests we made. For instance,
everyone knows the problem of checking the subscripts in array computation:
Algol-Compilers made (generally) a run-time checking of the computed subscript
expressions. This slows down the computation, but on the other hand helps for
avoiding errors. Nevertheless, that may discourage a programmer to prove the
programs he is writing. For instance, if you write a procedure to invert a
matrix, and you prove (formally or informally) that this procedure outputs the
correct inverse matrix, you may consider that the subscripts' run-time checks
are a waste of time. The only check you need is the checking of the squareness
of the input matrix. In a holon program, we will have a check of the matrix
shape at the ``invert matrix'' level, without checks of subscripts in the inner
levels. [This last point implies that there must exist functions to test the
components of a data structure as ``finding the bounds of the range of an array
subscript,'' etc\dots]

\bigskip

\noindent{\sl Remark}: The main point with the holon's tendency to autonomy is
that the operations check (at some level) the correctness of the type of the
data upon which they are operating. The following example appears in Koestler's
book: when a bird is building a nest, it follows a ``program'' made up of a set
of well-defined operations; it is not only by applying the operations in the
right sequence but also by applying the right operation sequence to the right
materials that the nest is built. For instance, the search and use of hard
material (twigs) must precede the catch and use of soft material (wool, moss).
For people with a structured programming mind, this description of a nest
building will look like a structured program!

%************************************************
\chapter{The Holon Language Syntax}
%************************************************

To define the language syntax we will use the BNF
notation.\marginnote{Transcriber's note: the author used an idiosyncratic
variation of Extended Backus-Naur form. The grammar that appears here is
written using the convention that the superscript $^+$ denotes the appearance
of one or more syntactic elements and `[\,\dots]' delimits optional elements.}

% ; nevertheless, we use the following symbols with the following meanings:

% 1) [ ]: these symbols delimit a syntactic element, the appearance of which is
% optional;

% 2) \{ \}: these symbols delimit a set of syntactic elements, one of which
% must be picked up;

% 3) The three dots \dots means that the previous symbol appears one or more
% times in the syntactic expression.

\section{Three kinds of syntax}

We will describe three kinds of syntax:

1) The external syntax of the language,

2) The internal syntax,

3) The modified external syntax.

\smallskip

\noindent The reader may puzzle about the third syntax; as a matter of fact
this syntax will look like the first one but with some abbreviations to ease
the handwriting (and in this last word, the prefix ``hand'' is important). The
modified syntax is a mapping of the external one in function of the kind of
tools we use to write the program.

\section{The external syntax}

\subsection{Programs, holons and names}

\begin{tcolorbox}
\sym{program} \is \sym{holon}\\
\sym{holon} \is \sym{holon name}\sym{holon body}\\
\sym{holon name} \is \sym{plain holon name}\\
    \qquad \alt \sym{holon name with parameters}\\
\sym{plain holon name} \is \sym{element of the unreserved symbols set}$^+$\\
\sym{holon name with parameters} \is\\
\qquad\sym{element of the unreserved symbols set}$^+$\\
\qquad\alt$($\sym{left identifier bracket}\sym{identifier}\sym{right identifier
bracket}$)^+$\\
\sym{element of the unreserved symbols set} \is\\
\qquad``any symbol except a reserved symbol''
\end{tcolorbox}

This last definition is not a BNF one, we used this informal sentence to
shorten the definition; as a matter of fact, the symbols set is not defined in
holon programming; you may use anything (ideograms if you will), but there are
some symbols which are reserved.

In the symbols set you use, you must reserve the following symbols:
\begin{enumerate}
    \item the identifier bracketing symbols (we use `\%' for left and right
identifier brackets),
    \item the connective symbol (we use `{\tt;}'),
    \item the reserved words used in the language statements (we set reserved
symbols in {\bf bold} font).
\end{enumerate}

The following examples are four valid holon names; the last one is not a
statement, it is a holon name.
\begin{itemize}
    \item Invert matrix \%A\% with bounds \%N\% \%M\%
    \item Search a free record
    \item Is the printer ready?
    \item $X := X+1$
\end{itemize}

\noindent{\sl Remarks}
\begin{enumerate}
    \item \sym{identifier} will be defined inside the data structures section.
    \item Two holon names are said to be ``identical'' if the symbols of the
two names match after the following operations:
    \begin{enumerate}
        \item Leading and trailing blanks are discarded,
        \item Strings of consecutive blanks are replaced by one blank symbol,
        \item The symbols between a left and its corresponding right identifier
brackets are ignored.

    The following holon names are ``identical'':
    \begin{itemize}
        \item Invert matrix \%X\% result in \%Y\%
        \item Invert matrix \%bobby\% result in \%ok\%
    \end{itemize}
    The following are not:
    \begin{itemize}
        \item Invert matrix \%A\%
        \item Invertmatrix \%4\%
    \end{itemize}
    \end{enumerate}
    \item There is not a limit to the length of a holon name; the holon name
may become a full explanatory text.
\end{enumerate}
\clearpage
\subsection{Holon body}

\begin{tcolorbox}
\sym{holon body} \is $[$\sym{datatest block}$]$$[$\sym{invariant
block}$]$\sym{function block}\\
\sym{function block} \is \sym{holon statements block}\\
\qquad\alt\sym{holon terminal language statements block}
\end{tcolorbox}

The important fact of the holon body syntax is the alternative for the writing
of the \sym{function block}: either like a \sym{holon statements block} or like
a \sym{holon terminal language statements block}; inside a holon there is no
mixing \sym{holon statement} (see further definitions) and \sym{holon terminal
language statement} (see the external syntax).

\subsection{Boolean holon}
\label{sec:boolean-holon}

If a holon name appears at well-defined positions in holon statements (these
positions will be defined in the holon statements syntax), this holon is called
a ``boolean holon.''

A boolean holon must contain at least one boolean expression (written in a
terminal language); the value ({\sl true} or {\sl false}) of this evaluated
expression is the value of the boolean holon.

If there is more than one boolean expression, the boolean holon value is the
value of the last expression.

It will be seen in a following section that the boolean holon may have access
to variables by value.

It must be clearly understood that a holon becomes a boolean holon by the
context; nevertheless a holon is either a ``pure'' holon or boolean one. If a
holon name appears in a spot such that it becomes a boolean one, all
appearances of this name in places where a boolean holon is not expected is an
error.

\subsection{Datatest and invariant blocks}

\begin{tcolorbox}
\sym{datatest block} \is \\
\qquad\kw{datatest} $($\sym{simple test} \alt \sym{test and error
recovery}$)^+$ \kw{end}\\
\sym{simple test} \is \sym{boolean holon name}{\tt;}\\
\sym{test and error recovery} \is \\
\qquad\sym{boolean holon name} \kw{else} \sym{holon name}{\tt;}\\
\sym{invariant block} \is \kw{invariant} \sym{simple test}$^+$ \kw{end}
\end{tcolorbox}

A data test evaluation is performed according to the following rules: [These
rules must be executed in this exact sequence]

\begin{enumerate}
    \item The \sym{test and error recovery} elements are evaluated. For each
element, the boolean holon is evaluated; if its value is {\sl false} then the
holon, the name of which follows the \kw{else} is executed.
    \item If any ``\kw{else} part'' was executed during the Phase 1, we
re-evaluate all the boolean holons appearing in the \sym{test and error
recovery} elements.
    \item The boolean holons of the \sym{simple test} elements are evaluated.
    \item If one (or more) {\it false} values appear during 2 and 3, the holon
execution is inhibited and a full report of the data test failure is printed
out.
\end{enumerate}

An invariant evaluation is performed according to the number 3 and number 4
operations of the data test evaluation rule. A data test is evaluated in front
of its holon execution, an invariant is evaluated inside a holon execution,
each time a connective symbol is met.

\bigskip

\noindent{\sl Remark}: The syntax doesn't require that the \sym{simple test}
elements are written after the \sym{test and error recovery} ones. The
sequencing---according to the evaluation rules---must be made by the compiler.
The programmer must be aware of that and be careful if you use some kind of
side-effect in the datatest.

\subsection{Holon statements block}

\begin{tcolorbox}
\sym{holon statements block} \is
\kw{begin} \sym{holon statement}$^+$ \kw{end}\\
\sym{holon statement} \is \sym{sequence}
\alt \sym{ifthenelse}
\alt \sym{dowhile}
\alt \sym{repeat}\\
\qquad\alt \sym{case}
\alt \sym{abort function}
\alt \sym{data structure commands}\\
\sym{sequence} \is $($\sym{holon name}{\tt;}$)^+$\sym{sequence unit}\\
\sym{sequence unit} \is \sym{holon name} \alt \sym{abort function}\\
\qquad\alt\sym{holon name}\sym{abort function}\\
\sym{ifthenelse} \is \kw{if} \sym{boolean holon name} \\
\qquad\qquad\qquad\qquad\kw{then} \sym{sequence unit}\\
\qquad\qquad\qquad\qquad$[$\kw{else} \sym{sequence unit}$]$\\
\sym{dowhile} \is \kw{do} \sym{holon name} \kw{while} \sym{boolean
holon}{\tt;}\\
\qquad\alt \kw{while} \sym{boolean holon} \kw{do} \sym{holon name}{\tt;}\\
\sym{repeat} \is \kw{repeat} \sym{holon name} \kw{until} \sym{boolean
holon}{\tt;}\\
\sym{case} \is \kw{case} $($\sym{boolean holon name} \kw{then}
\sym{sequence unit} \kw{eor}$)^+$\\
\qquad\qquad\qquad\qquad\qquad\qquad\qquad\quad\ \phantom{\kw{case} }\kw{else}
\sym{sequence unit}{\tt;}\\
\sym{abort function} \is \kw{exit} \alt \kw{stop}
\end{tcolorbox}

(The \sym{data structure commands} will be described in a following section
devoted to the data structures.)

In the above syntax, any appearance of the \sym{holon name} means that the
named holon is executed; any appearance of a \sym{boolean holon name} means
that the named boolean holon is evaluated.

We shall not describe thoroughly the semantics of the holon statements, because
the meanings of statements like \sym{ifthenelse}, \sym{dowhile}, \sym{repeat},
are well known.

An important point is: how do we look at the reserved words?

{\emergencystretch=1em The reserved words may be understood as ``flow of
control switches''; that is
the usual interpretation, enforced by the usual flowchart representation of a
program. Another interpretation is looking to the statements as functional
units and expressing their meanings using set
theory.\autocite{mills1972mathematical}\par}

The semantics of the \kw{case} are similar to the semantics of the conditional
form in LISP. The named boolean holons are evaluated in sequence. The first
boolean holon, the value of which is {\it true}, implies the execution of its
associated \sym{sequence unit}. This last execution completes the \kw{case}
statement. If no boolean holon is {\sl true}, then the last sequence unit is
executed (the one which follows the \kw{else}).

Introducing an \kw{exit} (or \kw{exitloop}) statement in the programming
languages isn't presently well-understood. As a matter of fact, nobody has
given a convincing explanation for its introduction or its rejection. We only
know that it is possible to write programs without any \kw{exit}, but that the
\kw{exit} eases the programming for some class of algorithms.

We have adopted the following policy: to admit two ``functions'': the \kw{exit}
and the \kw{stop}. The \kw{exit} means a local abortion of a holon
execution. The \kw{stop} implies a complete interruption of a program.

The syntax displays the limited use of the \kw{exit}. It may only appear in an
\kw{if} or a \kw{case} statement. The \kw{exit} may be viewed as a \kw{goto}
the last semicolon of a holon body.

In a loop, the \kw{exit} doesn't work like an ``{\sl exitloop}'' statement.
The syntax implies that the ``repeating part'' of a loop is abstracted in a
holon name. The ``condition part'' must be abstracted in a boolean holon. If in
the holon, the name of which is the ``repeating part,'' we use an \kw{exit},
this \kw{exit} will bring us back to the test of the condition part. An
example will be described in the appendix.

\subsection{Holon terminal language statements block}

\begin{tcolorbox}
\sym{holon terminal language statements block} \is \\
\qquad$($\sym{terminal language designator}\sym{terminal body}$)^+$\\
\sym{terminal body} \is \\
\qquad$($\sym{terminal language statement} \alt \sym{data structure
commands}$)^+$\\
\qquad\qquad\sym{end of terminal body}
\end{tcolorbox}

The three undefined elements of this section will be defined in the section
devoted to the ``Internal Syntax.''

\section{The internal syntax}

The external syntax describes how a holon is made of other holons. This process
stops when we meet a holon written in one of the terminal languages. The
terminal languages compose the internal syntax.

\subsection{Holon terminal languages}
\label{sec:holon-terminal-languages}

These languages contain only operative statements, the meaning of which may be
generalized by the following expression:

$$X \gets \mathrm{apply}(OP_1,OP_2,OP_3,\ldots ,OP_n)\ \mathrm{to}\
(Y_1,Y_2,Y_2,\ldots,Y_m)$$

In this expression: $X$ and $Y_i$ are the names of data structures; $OP_i$ is
the symbol of an operation; ``apply'' is a function, the rules of which are
well-defined, this function explains the order of the application of operations
on the data structures.

The expression can be read like this: ``the data structure $X$ takes a new
value. This value is computed by the application of a set of operations
$(OP_1,OP_2,OP_3,\ldots,OP_n)$ on the values of data structures
$(Y_1,Y_2,Y_3,\ldots,Y_m)$. The order of the operations application is defined
by the function `apply'.''

In this definition, data structure and operation must be understood in a wide
extent: a line on a printer is a data structure, and the term ``operation''
includes also functions like  sine, tangent, inverse, rewind, backspace, \dots
.

Only one thing is forbidden: a control statement. A terminal holon is a
sequence of expressions to be evaluated. These expression evaluations produce
results, but a terminal holon never makes a choice.

The expression previously written down isn't a syntax description. It is an
abstraction of the general meaning of instructions written in a lot of
languages. For instance, this expression abstracts all the following ones:
\begin{itemize}
    \item[] $X \gets A*B+C$
    \item[] print($Y*Z$)
    \item[] $Z = \sin ( \cos(2)) $
    \item[] backspace
    \item[] \kw{for} $X \gets 1$ \kw{step} $1$ \kw{until} $100$ \kw{do}
$A[X]\gets 0.0;$
\end{itemize}

\noindent (The last one isn't a control statement, it is a basic operation for
an array data structure.)

\medskip

\noindent We can imagine other expressions belonging to the abstract one:

\medskip

List children of Node Tree (A.B.C)
\begin{itemize}
    \item[] sort$_\uparrow$(B) \quad ($\uparrow$ meaning ascending order)
    \item[] sort$_\downarrow$(B) \quad ($\downarrow$ meaning descending order)
\end{itemize}

To sum up: the internal syntax is made up of a set of so-called ``terminal
languages,'' the nature of which depends on the program application areas.
These languages contain only operative statements, and no control statement.

The kinds and number of terminal languages will change from one user to
another, according to their specific interests. We will have terminal languages
devoted, for instance, to:
\begin{itemize}
    \item Integer Arithmetic
    \item Real Arithmetic
    \item Array Manipulation
    \item Character String Manipulation
    \item Text Edition
    \item Input-Output
    \item Graph Traversal
    \item Etc\dots
\end{itemize}

The holon system will include a list of ``standard'' terminal languages, plus
some ``special application oriented'' terminal languages. A mechanical engineer
would enjoy a ``gear oriented'' terminal language, etc\dots

Given the peculiar nature of the terminal languages, it is intended to define a
``generating holon'' which allows the programmer to define his own terminal
languages. The ``generating holon'' will produce a small compiler for each
defined terminal language.

\subsection{Defining a terminal language}

To define a terminal language, we must define the following elements:
\begin{enumerate}
    \item A ``name'' (the designator): it is a symbol used to point out this
terminal language among the available terminal languages. A terminal language
name designates one and only one terminal language (for the same computer
environment).
    \item A set of ``terminal data structure types'': i.e., a finite list of
names with correspondences to machine representation of the concept abstracted
by the ``type name.'' For instance: {\bf Integer}: 4 bytes.
    \item A set of ``terminal language operators'': i.e., a finite list of
symbols which abstract some manipulations and alterations upon data
structure(s). A description of the kinds of data structures upon which the
operator applies and also a description of the operator application results
must be furnished.
    \item A description of the ``apply'' function for the language: this
description explains the meaning of complex expressions involving many
operators and data structure names. (For instance, precedence order rules.)
This description implies also the definition of an ``end of statement'' symbol
\end{enumerate}

It must be understood that a sensible design of a terminal language includes
the definition of ``data structure properties test functions.'' That means
that, in order to write a powerful datatest, we must have the ability to test.
For instance, the definition of operations: ``lowbound'' and ``upbound'' which
we may use to determine the values $(n,m)$ for any array.

\subsection{Some syntax definitions}

We return to the undefined syntax elements appearing in the ``Statements
block.''

\begin{tcolorbox}
\sym{terminal language designator} \is \sym{most used language designator}\\
\qquad\alt \sym{special language designator}\\
\sym{most used language designator} \is {\tt\#}\\
\sym{special language designator} \is {\tt\#}\sym{string}\\
\sym{string} \is \sym{alphanu}\sym{space} \alt
\sym{alphanu}\sym{alphanu}\sym{space}\\
\qquad\alt \sym{alphanu}\sym{alphanu}\sym{alphanu}\sym{space}\\
\sym{alphanu} \is ``an alphabet character or a digit''\\
\sym{space} \is ``space or blank''
\end{tcolorbox}

(The last two syntactic definitions aren't BNF for the purpose of shortening
the exposition.)

    Examples of special language designator:
    \begin{itemize}
        \item[] \#GPH \quad standing for a graph language
        \item[] \#SUR \quad standing for a surveying language
        \item[] \#GR \quad standing for a gear computations language
        \item[] \#ARR \quad standing for an array alteration language
        \item[] Etc\dots
    \end{itemize}

    What we call the ``most used language'' is in the mind of the users the
terminal language which would be used the most frequently in terminal holons.
This language isn't necessarily the same at all places using the Holon
Programming Language.


\begin{tcolorbox}
\sym{terminal language statement}
\end{tcolorbox}

We don't describe this part of the syntax, we assume that the % reader can
% understand from section 5.3.1 what a terminal language syntax must look like.
reader can understand from the \nameref{sec:holon-terminal-languages} section
what a terminal language syntax must look like.

\begin{tcolorbox}
\sym{end of terminal body} \is {\tt\#\#}
\end{tcolorbox}

\section{Extending the character set}

This section will describe a feasible extension of the concept of terminal
languages designed by the programmer. This extension was suggested to me by
reading some of Babbage's Paper published by
Morrison.\autocite{morrison1961charles}

\subsection{Babbage's language for technical drawing}

Charles Babbage designed and used a language of his own to explain the working
of the mechanical devices embedded in his calculating engines. What struck me
reading Babbage's Papers was his clear distinction between a machine's static
description (blueprint) and a machine's working description (statement in his
language). Joining a machine's working formal description to a blueprint isn't
usual.

Babbage explained that the need of this language appeared in proportion as the
design of his engine progressed. It was the growing complexity of the
calculating engine design which implies the need of the language describing the
working of the machine elements. This may be viewed as an application of the
principles expounded by Simon and Koestler.

They claim the need for a neat distinction between the world as sensed and the
world as acted. This distinction ensures correct development of complex
systems.

Now, looking to Babbage's Language, you may infer that with the help of a
computer, it is feasible to do a lot of things with this syntax; for instance:
\begin{enumerate}
    \item To use the working formal description for driving a moving display of
the mechanical parts on a CRT-Device.
    \item To expand a ``macro'' description of the working to the complete
formal one.
\end{enumerate}
But any attempt to implement systems like these is a still-born one. Why?
Because the usefulness of Babbage's syntax depends on the use of a new set of
small symbols, the reading of which can become as meaningful for a mechanical
engineer as looking at notes of a score can be for a musician. The mechanical
engineering oriented reader is urged to read Babbage's paper {\it On a method
of Expressing by Signs the Action of Machinery}.\footnote{Ibid.}

Considering that the previous example isn't unique, (for instance: think of the
APL character set), our statement is: in many applications, computer-aided
systems are cumbersome and awkward because the methods to extend the character
set are cumbersome and awkward.

People can presently extend the character set only by the use of light-pen (for
the input) and plotter or CRT (for the output). It is also possible to build
character sets dedicated to an I/O device like the APL terminal, but this
method is obviously unflexible.

In the next section, we describe a flexible device to input ``personal''
character sets.

\subsection{A symbol generator}

The device we describe will be a new combination of existing components. In
spite of my lack of electronic knowledge, I think that ``wiring'' this device
wouldn't be a nightmare.

To represent a symbol, we can use a pattern of dots. This Dots-Pattern method
is already extensively used by many printing devices. We don't care about the
way (impact, ink spitting, thermal, electrostatic\dots) we use to lay down the
expected dots-pattern on the paper. We don't care anymore about the pattern
shape (number of dots); we must have enough dots (perhaps $8\times8$) for
generating fair copy of ``continuously written'' symbols.

Consider a keyboard composed of three sections:
\begin{enumerate}
    \item A standard typewriter (or teletype) keyboard;
    \item A set of keys with ``undefined'' characters;
    \item A set of keys to manage the assignment of characters to keys of the
previous set.
\end{enumerate}
Each time we depress a key of the standard keyboard, a dots-pattern of the
corresponding character appears somewhere (printed on a paper or lighted by LED
(Light Emitting Diode)) and we assume that the binary-code of this character is
recorded somewhere.

When we depress a key of the ``undefined'' character set, we can record a
binary-code corresponding to the depressed key. This binary-code is independent
of any dots-pattern. When this binary-code appears on an input or output
device, this appearance can start a search for the corresponding dots-pattern
in memory. When this dots-pattern is found in memory, its description is used
to drive the dot-writing device.

The user depressed keys in the third key set to build up the content of the
memory. For the input device, the support of the memory would be something like
a magnetic card; each user introduces his own card upon which he can record his
own character set.

When a user puts a recorded card into the device, the ``undefined'' keys will
take ``values'' according to the memory card content. That can be displayed to
the user by fitting into the hollow of each ``undefined'' key a matrix of LED
Dots which will be turned on according to the memory card content.

The keys of the third set are used mainly to point out the coordinates of the
dots which must be ``activated'' to form the symbol. The lighted dots would
appear first on an enlarged dots-display to help the user to draw the symbol.
Output devices will need memory to manage the mapping of the ``undefined''
characters into dot patterns. In the output character string, there will exist
control commands to write content inside the memory of the output device.
Commands such as: ``The next `X' bytes in this output string must write inside
the memory.''

\subsection{Use in terminal language}

The programmer defining a new terminal language can build a symbolic syntax
that ``tightens'' the requirements of the application area. The new symbols can
stand for operators or data type constants. (This use can simplify some input
procedure, for instance in hotel reservation systems, the wishes of the
customer can be input like the conventional ideograms used by this business.)

The set of the new symbols with their meanings---in terms of holons
(procedures) and data values---will be a part of the terminal language
definition.

\section{Inner and outer syntaxes}

The Holon statements are used to write the several program levels, all but the
terminal ones where the programmer uses the terminal languages. The programmer
can define his personal terminal languages, provided he fulfills the
requirement that his new syntax corresponds only to an ``operator grammar.''

A program is completed when all the levels are expressed in terms of terminal
languages.

As a matter of fact, there exists a clear partition between an outer syntax and
several inner syntaxes.\autocite{wilkes1968outer}

\section{Data structure commands and modified syntax}

This syntax description must be completed by the description of the ``Data
Structure Commands'' and also of the ``Modified Syntax.'' However, we delay
these descriptions. The description of the semantics of these syntax elements
will be more simple after the reading of the next section devoted to the
representation of structured programs.

%************************************************
\chapter{Structured Program Representation}
%************************************************

A structured program may be represented by
flowcharting,\autocite{mills1972mathematical} but this representation does not
visualize the several levels if we embed (in the flowchart) the levels in each
other. A tree representation seems more
adequate.\autocite{henderson1972experiment} We have chosen a binary tree
representation for a holon program.\autocite{knuth1968art} This representation
is used as a model of a program execution and also as a base for the
implementation of the holon concept.

\section{Binary tree representation}

Figure~\ref{fig:tree} depicts a tree representation of a structured program (on
the left) with its corresponding corresponding binary tree (on the right).

\begin{figure}
\begin{center}
% Figure 1
\begin{forest}
for tree={
% style of nodes in the tree
draw, semithick, circle,
       minimum height = 3ex,
        minimum width = 2em,
% style of tree (edges, distances, direction)
     grow = south,
%forked edge,        % for forked edge
    s sep = 3mm,    % "sibling" distance
    l sep = 3mm,    % "level" distance
 fork sep = 3mm,    % distance from parent to branching point
}
[A,fit=band
[B [E][F]]
[C [G]]
[D [H][I]]
]
\end{forest}
\qquad
% Figure 2
\begin{tikzpicture}
  [vertex/.style={circle,semithick,draw=black,fill=white},
   node distance=2cm, minimum size=20pt,
   on grid,
   >=latex,>={Stealth[inset=0pt,length=6pt,angle'=28,round]}
  ]
  \node[vertex] (A) {A};
  \node[vertex,below left=1cm and 0.75cm of A] (B) {B};
  \node[vertex,below left=1cm and 0.75cm of B] (E) {E};
  \node[vertex,right=2cm of B] (C){C};
  \node[vertex, right=1cm of E] (F){F};
  \node[vertex, below left=1cm and 0.75cm of C] (G){G};
  \node[vertex,right=1cm of C] (D){D};
  \node[vertex, below left=1cm and 0.75cm of D] (H){H};
  \node[vertex, right=1cm of H] (I) {I};
  \draw[->]
    (A) edge (B)
    (B) edge (E)
    (B) edge (C)
    (E) edge (F)
    (C) edge (G)
    (C) edge (D)
    (D) edge (H)
    (H) edge (I);
\end{tikzpicture}
\end{center}
\caption{A structured program's tree and binary tree representations.}
\label{fig:tree}
\end{figure}

The contents of the nodes of these trees must be clearly understood: A node
contains only a holon name, unless the holon is a terminal one (written in
terminal language); the node corresponding to a terminal holon contains the
holon name and the terminal instructions.

The holon tree of Figure~\ref{fig:tree} corresponds to the syntactic form:
\begin{itemize}
    \item[A] \kw{begin} B; C; D; \kw{end}
    \item[B] \kw{begin} E; F; \kw{end}
    \item[C] \kw{begin} G; \kw{end}
    \item[D] \kw{begin} H; I; \kw{end}
\end{itemize}

[We temporarily use the ``sequence'' statement only.] Thus, for the binary tree
we have the following rules:
\begin{enumerate}
    \item A node contains a holon name of a terminal holon;
    \item The left link points to the first statement of the holon body;
    \item The right link points to the next statement (if any) in the holon
body.
\end{enumerate}

\noindent [It can be said that a left link is a graphical correspondence of
\kw{begin} and the right link is the one of a semicolon which is the
connective symbol.]

\section{Tree traversal and program execution}

A program execution may be viewed as a tree traversal with respect to some
order. The order must be chosen so that the execution of the terminal level
instructions is performed in the correct order. For the binary tree, the
traversal will be the so-called postorder one (or symmetric one).

This traversal is recursively defined by Knuth:\footnote{Ibid.}

\medskip

\noindent\-\hspace{0.4cm} \kw{if} the tree is empty\\
\-\hspace{0.7cm} \kw{then} do nothing\\
\-\hspace{0.7cm} \kw{else} \kw{begin} \\
\-\hspace{0.9cm} traverse the left subtree;\\
\-\hspace{0.9cm} visit the root;\\
\-\hspace{0.9cm} traverse the right subtree;\\
\-\hspace{0.7cm} \kw{end}

\medskip

We use the following ``refinement'' for the semantic level (holon) ``visit the
root'':

\medskip

\noindent\-\hspace{0.4cm} \kw{if} the root is a terminal holon\\
\-\hspace{0.7cm} \kw{then} perform the contained instructions\\
\-\hspace{0.7cm} \kw{else} do nothing\\

\medskip

If in Figure~\ref{fig:tree}, we assume that nodes E,F,G,H, and I are terminal
holons, a postorder traversal will start the terminal holon executions in the
correct order.

Of course, this binary tree model works only for the sequence statement of
Chapter 5's syntax; but we shall develop the model in order to include the
other holon statements.
Nevertheless, the expansion of the model must respect the clear graphical
representation of the levels that the binary tree allows.

\section{Binary tree representation of the holon statements}


\noindent We shall imagine that there exists a ``textual
index''\autocite{dijkstra1968letters} which is driven by the postorder
traversal algorithm. When the ``textual index'' enters in a terminal holon, it
starts the terminal holon statements.

To simulate repetitive statements and conditional statements, we include in the
binary tree some ``switching'' nodes which will be ``switched'' by boolean
holon values. These nodes will force the ``textual index'' to re-traverse some
subtrees or will block the way to some subtrees, according to boolean holon
values.

[The difference between this model and the flowchart will be the clear
distinction of the levels.]

The model will be composed from the following types of nodes:
\begin{marginfigure}[5.05cm]
    % Figure sidebar
\begin{tikzpicture}%
  [vertex/.style={circle,semithick,draw=black,fill=white},
  square/.style={regular polygon,regular polygon
sides=4,semithick,draw=black,fill=white,minimum size=14pt},
   node distance=2cm, minimum size=14pt,
   on grid,
   >=latex,>={Stealth[inset=0pt,length=6pt,angle'=28,round]}
  ]
  \node[vertex] (O) {};
  \node[vertex, below=0.8cm of O] (A) {A};
  \node[square,below=1cm of A] (C) {\footnotesize C};
  \node[square,below=0.7cm of C] (F) {\footnotesize F};
  \node[above right=0.2cm and 0.49cm of F] {\scriptsize DO};
  \node[square,below=0.7cm of F] (G) {\footnotesize G};
  \node[above right=0.2cm and 0.48cm of G] {\scriptsize RE};
  \node[vertex, below=0.95cm of G] (d) {d};
  \node[vertex, below=0.9cm of d] (r) {r};
\end{tikzpicture}
\end{marginfigure}
\begin{itemize}
    \item[a] The ``empty node,'' used to fix levelling inside the tree;
    \item[b] The ``holon name node,'' representing a holon name, (The depicted
symbol corresponds to a holon, the name of which is: \textsf{A});
    \item[c] The ``boolean colon name node,'' representing a boolean holon
name; there are three distinct symbols according to the statement where the
boolean holon name appears
    \setlength\itemsep{1.8em}
    \item[d] The ``dispatching node,'' used to implement a choice;
    \setlength\itemsep{.8em}
    \item[e] The ``repeating node,'' used to implement a repetition.
\end{itemize}

\bigskip
\subsection{if then else}

\medskip
\noindent Let the holon statement be:

\medskip

\noindent\-\hspace{0.7cm}\kw{if} H \kw{then} A \kw{else} B.

\medskip
% Figure 2
\begin{marginfigure}[1cm]
\begin{tikzpicture}
  [vertex/.style={circle,semithick,draw=black,fill=white},
  square/.style={regular polygon,regular polygon
sides=4,semithick,draw=black,fill=white,minimum size=26pt},
   node distance=2cm, minimum size=20pt,
   on grid,
   >=latex,>={Stealth[inset=0pt,length=6pt,angle'=28,round]}
  ]
  \node[vertex] (O) {};
  \node[square,below left=1cm and 0.25cm of O] (H) {H};
  \node[vertex,right=1.5cm of H] (d){d};
  \node[vertex, below left=1cm and 0.75cm of d] (A){A};
  \node[vertex, below right=1cm and 0.75cm of d] (B){B};
  \draw[->]
    (O) edge (H.north)
    (H) edge (d)
    (d) edge (A)
    (d) edge (B);
\end{tikzpicture}
\caption{Binary tree model for an \mbox{if-then-else} statement.}
\label{fig:ifthenelse}
\end{marginfigure}

\noindent The binary tree model is depicted in Figure~\ref{fig:ifthenelse}.

In Figure~\ref{fig:ifthenelse}, the holons defining the boolean holon H and the
holons A and B must be linked to the corresponding nodes by left linked
subtrees. H, A, and B are not necessarily terminal holons.

A postorder traversal of the tree in Figure~\ref{fig:ifthenelse} implies the
following order for the visited nodes: H, A, d, B, ``empty node.''

(When we say that a node is visited, that means that the left subtree rooted by
this node has been traversed. During this traversal, the terminal nodes are
executed in the correct order.)

The visit of H gives a value to the boolean holon H (unless there is a
programming error). The value is used to modify the content of the dispatching
node links.

A dispatching node may be viewed (see Figure~\ref{fig:dispatch}) as a node with
four links:
\begin{itemize}
    \item Two structural ones which can never be modified during program
execution
    \item Two operational ones which are modified by the boolean holon value.
\end{itemize}

% Figure 3
\begin{marginfigure}
\begin{tikzpicture}%
  [vertex/.style={circle,thick,draw=black,fill=white},
   hole/.style={draw,shape=circle, semithick, fill=white,minimum size=0.15cm,
inner sep=0pt},
   node distance=2cm,
   on grid,
   >=latex,>={Stealth[inset=0pt,length=6pt,angle'=28,round]}
  ]
  \node[vertex, thick, minimum size=50pt] (D) {};
  \node[hole, above left=0.35cm and 0.35cm of D] (NW) {};
  \node[left=.9cm of NW] (dummywest) {};
  \node[hole, below left=0.35cm and 0.35cm of D] (SW) {};
  \node[hole, above right=0.35cm and 0.35cm of D] (NE) {};
  \node[right=.9cm of NE] (dummyeast) {};
  \node[hole, below right=0.35cm and 0.35cm of D] (SE) {};
  \node[vertex, below left=1.9cm and 1.4cm of D] (A) {A};
  \node[vertex, below right=1.9cm and 1.4cm of D] (B) {B};
  \node[above right=0.8cm and 0.8cm of D] () {d};

  \node[right=1.3cm of SE] (opkeyleft) {};
  \node[right=2.5cm of SE] (opkeyright) {};
  \node[right=0.8cm of opkeyright] {\footnotesize operational link};
  \draw[semithick,dashed,->] (opkeyleft)--(opkeyright);

  \node[below=0.5cm of opkeyleft] (stkeyleft) {};
  \node[below=0.5cm of opkeyright] (stkeyright) {};
  \node[right=0.7cm of stkeyright] {\footnotesize structural link};
  \draw[semithick,->] (stkeyleft)--(stkeyright);

  \draw[semithick] (NW)--(dummywest.center);
  \draw[semithick] (NE)--(dummyeast.center);
  \draw[semithick,->] (dummywest.center)--(A.north);
  \draw[semithick,->] (dummyeast.center)--(B.north);
  \draw
    (D.south)--(D.north)
    (D.west)--(D.east);
  \draw[dashed,semithick,->]
    (SW) edge (A)
    (SE) edge (B);
\end{tikzpicture}
\caption{A dispatching node.}
\label{fig:dispatch}
\end{marginfigure}

One of the operational links is ``broken'' according to the boolean holon
value.

``visit the node'' for a boolean holon will look like:

\smallskip

\noindent\-\hspace{0.4cm}\kw{datatest} has a value been given to the boolean
holon?
\kw{end} \\
\noindent\-\hspace{0.4cm}\kw{begin}\\
\noindent\-\hspace{0.7cm}operational links $\gets$ structural links;\\
\noindent\-\hspace{0.7cm}\kw{if} boolean holon value $=$ {\it true}\\
\noindent\-\hspace{0.9cm}\kw{then} right operational link $\gets$ nil\\
\noindent\-\hspace{0.9cm}\kw{else} left operational link $\gets$ nil;\\
\noindent\-\hspace{0.7cm}erase boolean holon value;\\
\noindent\-\hspace{0.4cm}\kw{end}

\smallskip

The result of the previous operation will be the ``pruning'' of one of the
subtrees of the labelled -d tree.\marginnote[-0.3cm]{This refers to the
dispatch node, i.e., one of the branches is pruned.}

\subsection{do while}
\label{sec:do-while}

Let the holon statement:

\smallskip

\noindent\-\hspace{0.7cm}\kw{do} A \kw{while} C;

\smallskip

\noindent be modeled by the binary tree depicted in Figure~\ref{fig:dowhile}.

% Figure 4
\begin{marginfigure}
\begin{tikzpicture}%
  [vertex/.style={circle,semithick,draw=black,fill=white},
  square/.style={regular polygon,regular polygon
sides=4,semithick,draw=black,fill=white,minimum size=26pt},
   node distance=2cm, minimum size=20pt,
   on grid,
   >=latex,>={Stealth[inset=0pt,length=6pt,angle'=28,round]}
  ]
  \node[vertex] (O) {r};
  \node[square,below left=1.15cm and 0.75cm of O] (C) {C};
  \node[vertex,right=1.5cm of C] (d){d};
  \node[vertex, below left=1.15cm and 0.75cm of d] (A){A};
  \draw[->]
    (O) edge (C.north)
    (d) edge (A);
\draw[->] (C) -- (d) node [midway, above, xshift=-0.05cm, yshift=-0.15cm] {{\tt
DO}};
\end{tikzpicture}
\caption{Binary tree model for a \mbox{do-while} statement.}
\label{fig:dowhile}
\end{marginfigure}

In this case, we use a ``repeating'' node. This kind of node has two structural
links (the right one points to the tree mapping the holon next statement; it is
not depicted in Figure~\ref{fig:dowhile}) and one left operational link
modified by the boolean holon.

The symbol ``do'' on the side of the boolean holon is used to select the
adequate ``visit the node'' refinement.

This refinement will be:

\smallskip

\noindent\-\hspace{0.4cm}\kw{datatest} has a value been given to the boolean
holon? \kw{end}
\\
\noindent\-\hspace{0.4cm}\kw{begin} \\
\noindent\-\hspace{0.7cm}d-left operation link $\gets$ d-left structural
link;\\
\noindent\-\hspace{0.7cm}\kw{if} boolean holon value $\equiv$ {\it false}\\
\noindent\-\hspace{0.9cm}\kw{then} \\
\noindent\-\hspace{1.1cm}d-left operational link $\gets$ nil;\\
\noindent\-\hspace{1.1cm}r-left operational link $\gets$ nil;\\
\noindent\-\hspace{0.9cm}\kw{else} do nothing; \\
\noindent\-\hspace{0.7cm}erase boolean holon value;\\
\noindent\-\hspace{0.4cm}\kw{end}

\subsection{repeat until}

Let the holon statement:

\smallskip

\noindent\-\hspace{0.7cm}\kw{repeat} B \kw{until} Y;

\smallskip

\noindent correspond to the binary tree shown in Figure~\ref{fig:repeatuntil}

% Figure 5
\begin{marginfigure}
\begin{tikzpicture}%
  [vertex/.style={circle,semithick,draw=black,fill=white},
  square/.style={regular polygon,regular polygon
sides=4,semithick,draw=black,fill=white,minimum size=26pt},
   node distance=2cm, minimum size=20pt,
   on grid,
   >=latex,>={Stealth[inset=0pt,length=6pt,angle'=28,round]}
  ]
  \node[vertex] (r) {r};
  \node[vertex,below left=1cm and 0.75cm of r] (B) {B};
  \node[square,right=1.5cm of B] (Y){Y};
  \draw[->]
    (r) edge (B)
    (B) edge (Y);
\draw node [right=0.6cm of Y,yshift=0.25cm] {{\tt RE}};
\end{tikzpicture}
\caption{Binary tree model for a \mbox{repeat-until} statement.}
\label{fig:repeatuntil}
\end{marginfigure}

The refinement of the ``visit the node'' for the boolean holon:

\smallskip

\noindent\-\hspace{0.4cm}\kw{datatest} has a value been given to the boolean
holon? \kw{end}
\\
\noindent\-\hspace{0.4cm}\kw{begin} \\
\noindent\-\hspace{0.7cm}\kw{if} boolean holon value $\equiv$ {\it true}\\
\noindent\-\hspace{0.9cm}\kw{then} r-left operational link $\gets$ nil;\\
\noindent\-\hspace{0.7cm}erase boolean holon value;\\
\noindent\-\hspace{0.4cm}\kw{end}

\subsection{case statement}

Let the holon statement:

\smallskip

\noindent\-\hspace{0.4cm}\kw{case} M \kw{then} Z\\
\noindent\-\hspace{0.7cm}\kw{eor} N \kw{then} W\\
\noindent\-\hspace{0.7cm}\kw{eor} \kw{else} T;

\smallskip

\noindent correspond to the binary tree in Figure~\ref{fig:case}.

% Figure 6
\begin{marginfigure}
\begin{tikzpicture}%
  [vertex/.style={circle,semithick,draw=black,fill=white},
  square/.style={regular polygon,regular polygon
sides=4,semithick,draw=black,fill=white, minimum size=26pt},
   node distance=2cm, minimum size=20pt,
   on grid,
   >=latex,>={Stealth[inset=0pt,length=6pt,angle'=28,round]}
  ]
  \node[vertex] (o) {};
  \node[square,below left=1.25cm and 0.75cm of o] (M) {M};
  \node[vertex,right=1.2cm of M] (d){d};
  \node[square,below left=1.25cm and 0.75cm of d] (Z) {Z};
  \node[square,right=1.2cm of d] (N){N};
  \node[vertex,right=1.2cm of N] (d'){d};
  \node[square,below left=1.25cm and 0.75cm of d'] (W) {W};
  \node[vertex,right=1.2cm of d'] (T){T};
  \draw[->]
    (o) edge (M.north)
    (M) edge (d)
    (d) edge (N)
    (N) edge (d')
    (d') edge (T)
    (d) edge (Z.north)
    (d') edge (W.north);
\end{tikzpicture}
\caption{Binary tree model for a case statement.}
\label{fig:case}
\end{marginfigure}

The ``visit the node'' refinement used for the boolean holon in this statement
is similar to the one used for the {\it if-then-else} statement.

\section{Full holon representation}

A full holon contains a datatest, an invariant and a sequence of holon
statements. For instance:

\smallskip

\noindent\-\hspace{0.4cm}N\\
\noindent\-\hspace{0.4cm}\kw{datatest} \dots \kw{end} \\
\noindent\-\hspace{0.4cm}\kw{invariant} \dots \kw{end} \\
\noindent\-\hspace{0.4cm}\kw{begin} S$_1$; S$_2$; S$_3$; \kw{end} \\

\smallskip

The S$_i$ are holon statements (not only holon names). The binary tree is
displayed in Figure~\ref{fig:fullholon}.

\begin{figure}
% Figure 8
\begin{tikzpicture}%
  [vertex/.style={circle,semithick,draw=black,fill=white},
  ell/.style = {% E as ellipse
    ellipse,semithick,draw=black,fill=white},
   node distance=2cm, minimum size=20pt,
   on grid,
   >=latex,>={Stealth[inset=0pt,length=6pt,angle'=28,round]}
  ]
  \node[vertex] (N) {N};
  \node[ell,below left=1.25cm and 0.5cm of N] (D) {DATATEST};
  \node[below left=1.875cm and 0.75cm of D] (D0) {};
  \node[ell,right=2.5cm of D] (S1) {S$_1$};
  \node[below left=1.875cm and 0.75cm of S1] (S10) {};
  \node[ell,right=1.6cm of S1] (I1) {INV};
  \node[below left=1.875cm and 0.75cm of I1] (I10) {};
  \node[ell,right=1.6cm of I1] (S2) {S$_2$};
  \node[below left=1.875cm and 0.75cm of S2] (S20) {};
  \node[ell,right=1.6cm of S2] (I2) {INV};
  \node[below left=1.875cm and 0.75cm of I2] (I20) {};
  \node[ell,right=1.6cm of I2] (S3) {S$_3$};
  \node[below left=1.875cm and 0.75cm of S3] (S30) {};
  \node[ell,right=1.6cm of S3] (I3) {INV};
  \node[below left=1.875cm and 0.75cm of I3] (I30) {};
  \draw[->]
    (N) edge (D)
    (D) edge (S1)
    (S1) edge (I1)
    (I1) edge (S2)
    (S2) edge (I2)
    (I2) edge (S3)
    (S3) edge (I3);
  \draw[dashed,semithick]
    (D)--(D0)
    (S1)--(S10)
    (I1)--(I10)
    (S2)--(S20)
    (I2)--(I20)
    (S3)--(S30)
    (I3)--(I30);
\end{tikzpicture}
\caption{Binary tree model for a full holon.}
\label{fig:fullholon}
\end{figure}

The dashed lines are links to the defining holon trees for each of the holon
components.

The nodes $S_1$, $S_2$, $S_3$ are the roots of left subtrees corresponding to
holon statement trees. For instance: if $S_1$ is a holon name, the pending link
points to the definition of this holon in terms of datatest, invariant and
holon statements; if $S_i$ is a \kw{do while} statement, the node marked $S_i$
in Figure~\ref{fig:fullholon} is the ``repeating'' node which ``roots'' a {\it
dowhile} binary tree (see the~\nameref{sec:do-while} section); etc.

[The description of the datatest and invariant representation is given in one
of the appendices.]

\section{The postorder traversal algorithm revisited}

We can rewrite the refinement of the holon ``visit the node'':

\smallskip

\noindent\-\hspace{0.4cm}\kw{case} node is a terminal one\\
\noindent\-\hspace{0.7cm}\kw{then} execute the contained terminal language
statement\\
\noindent\-\hspace{0.4cm}\kw{eor} node is a repeating one with left operational
link $\neq$ nil\\
\noindent\-\hspace{0.7cm}\kw{then} traverse left subtree and revisit this
node\\
\noindent\-\hspace{0.4cm}\kw{eor} node is a repeating one (with left
operational
link$=$nil)\\
\noindent\-\hspace{0.7cm}\kw{then} left operational link $\gets$ left
structural
link\\
\noindent\-\hspace{0.4cm}\kw{eor} node is a boolean one\\
\noindent\-\hspace{0.7cm}\kw{then} operate according to the boolean type (do,
re,
empty)\\
\noindent\-\hspace{0.4cm}\kw{eor} \kw{else} do nothing.

\smallskip

It must be understood that this postorder traversal algorithm can be used to
``interpret'' the program. The textual index access to and exit from a holon
can be reported to display the progress and the history of the computation.

\bigskip

\noindent[{\it Remark}: In this section all the algorithms were described
according to the rules of the holon syntax.]

%************************************************
\chapter[Mapping the Binary Tree into an Executable Program]{Mapping the Binary
Tree into an\\Executable Program}
%************************************************

It must be clearly understood that the syntax of the holon programming
language is designed in order to write well-structured and well-documented
programs. A holon name can contain a full explanation of its abstracted
operation.

We shall explain in this section how this well-structured and well-documented
program will be mapped into a sequence of efficient machine-code instructions.
In this instruction sequence, the program structure will not appear explicitly.
But the code generation will be driven by the program structure.

Each holon will make use of its ``tendency to integration'' during the
code-generation operation.

\section{Basic integration}
\label{sec:basic-integration}

Let us consider a holon program with the following properties:
\begin{itemize}
    \item[a.] The program is complete: i.e., all the ``leaves'' (the pending
nodes) of its binary tree representation are holons written in terminal
languages.
    \item[b.] The binary tree representation of the program is a minimal one:
i.e., the definition of a holon appears only in one place in the tree; thus a
duplication of a holon name does not imply a copy of the subtree rooted by the
node with the same name.
    \item[c.] The binary tree representation of the program is a
left-definition one: i.e., the subtree defining a holon is ``attached'' to the
leftmost node labelled with this holon name; this means that during a postorder
traversal using the {\it structural} links, the textual index has access either
to the subtree defining a holon or to a holon, the defining subtree of which
has been already traversed.
\end{itemize}
[The traversal uses the structural links to avoid any subtree pruning or
retraversal.]

\noindent We suppose also that:
\begin{itemize}
    \item[a.] There exist algorithms to generate machine code for each terminal
language.
    \item[b.] We can assign storage addresses for the data structures, and
these storage addresses are used in the machine code instructions.
\end{itemize}
We have the choice to implement each holon: either as a procedure (closed
routine) or an inline inclusion (open routine). The choice can be based upon a
series of criteria:
\begin{enumerate}
    \item The length of the machine-code generated for the holon.
    \item The frequency of the holon name appearance in the program.
    \item The relative cost of a procedure call for the used computer.
\end{enumerate}
These three criteria are the expression of the computer architecture influence.
They are put together in some trade-off function.

Nevertheless, the choice resulting from the trade-off function may be modified
in some cases, according to the context of a holon name appearance in the
program. This leads to a new criterion:
\begin{itemize}
    \item[4.] The kind of instruction which surrounds a holon name. (For
instance, after a \kw{do} an open routine can be used instead of a routine
call.)
\end{itemize}
But this last criterion has only a local influence; i.e., the first three
criteria are used to make a global decision about the implementation unit for a
holon, but locally, the fourth criterion may modify the choice.

\section{Recursive holon}

The criteria of the previous section can be applied only to non-recursive
holons. If a holon is a recursive one, it needs a special implementation
mechanism.\autocite{dijkstra1960recursive} Before applying the criteria, we
must ``pin'' the recursive holons.

In the conventional languages, the operation is not an easy one. Generally, the
system is designed to avoid this problem either by assuming that all procedures
may be recursive (ALGOL-compiler) or by assuming that the programmer will
specify the recursive procedures (PL/I syntax), the latter choice assumes that
the programmer is able to specify direct- and cross-recursive procedures.

The only algorithm that I know to detect the recursive procedure is the
transitive closure one (see Wegner, section
4.7.6).\autocite{wegner1968programming} This algorithm is simple, but when you
read the description with an example, you may wonder how it can be manageable
if the number of procedures is large.

We shall show that if a program is a structured one, we can design a simple
algorithm to detect recursive holons. [The difference between this algorithm
and the transitive closure one is that it seems insensitive to the number of
procedures, with respect to deciding if a procedure is potentially recursive or
not.]

We suppose that the program has a binary tree representation which is a minimal
one (see~\nameref{sec:basic-integration}). We assume that we have a list of
each holon name with a list of its appearances inside the tree. Appearances are
pointed out by a so-called ``Dewey Binary Notation for Binary
Trees.''\autocite{knuth1968art} Detecting recursivity will become a simple
bit-strings comparison.

Suppose a holon A define by Figure~\ref{fig:dewey}.

\begin{marginfigure}[1cm]
% Figure 9
\begin{tikzpicture}%
  [vertex/.style={circle,semithick,draw=black,fill=white},,
   node distance=2cm, minimum size=20pt,
   on grid,
   >=latex,>={Stealth[inset=0pt,length=6pt,angle'=28,round]}
  ]
  \node[vertex] (A) {A};
  \node[vertex,below left=1.25cm and 0.75cm of A] (D) {D};
  \node[vertex,right=1.5cm of D] (B){B};
  \node[vertex,below left=1.25cm and 0.75cm of B] (D') {D};
  \node[vertex,right=1.5cm of D'] (C){C};
  \node[vertex,below left=1.25cm and 0.75cm of C] (A') {A};
  \draw[->]
    (A) edge (D)
    (D) edge (B)
    (B) edge (D')
    (D') edge (C)
    (C) edge (A');
\draw node[left=0.5cm of A,yshift=0.25cm] {{\tt 1}};
\draw node[left=0.6cm of D,yshift=0.25cm] {{\tt 10}};
\draw node[left=0.8cm of D',yshift=0.25cm] {{\tt 1010}};
\draw node[right=0.7cm of B,yshift=0.25cm] {{\tt 101}};
\draw node[right=0.9cm of C,yshift=0.25cm] {{\tt 10101}};
\draw node[right=1cm of A',yshift=0.25cm] {{\tt 101010}};
\end{tikzpicture}
\caption{A binary tree labeled with Dewey binary notation.}
\label{fig:dewey}
\end{marginfigure}

The holon names appearance list is:

\begin{itemize}
\setlength{\itemindent}{1cm}
\setlength\itemsep{-0.5em}
    \item[A:] 1, 101010
    \item[B:] 101
    \item[C:] 10101
    \item[D:] 10, 1010
\end{itemize}

Is ``D'' recursive?

``D'' appears at nodes 10 and 1010; bit-string comparison (starting from the
left) shows that the first dissimilar digit is a `1'. This means the second
appearance of D belongs to the right subtree of the first D. Thus, D is not
included in its definition (a definition is always a left subtree). Thus, D is
not recursive.

Is ``A'' recursive?

``A'' appears at nodes 1 and 101010; bit-string comparison shows that the first
dissimilar digit is a `0'. Thus the second A is included in the left subtree of
the first A. Thus, A is recursive. But this recursivity implies that B and C
are recursive, because in a postorder traversal, when the textual index has
access to the second appearance of A, the traversal of the B and C subtrees is
not complete.

How do we detect cross-recursivity (like B and C)?

We take the bit-string identification of the second appearance; for A, it is:
101010. And we extract all the substrings (starting at the left), the first
bits of which match the bit-string of the first appearance and the last bit of
which is followed by a `0'.

For the example, this procedure will produce the set of bit-strings:
\begin{itemize}
\setlength{\itemindent}{1cm}
\setlength\itemsep{-0.2em}
    \item[1]    which corresponds to A
    \item[101] which corresponds to B
    \item[10101] which corresponds to C
\end{itemize}
[This bit-string comparison is a back-tracking from a descendant to his
ancestors, in a family lineal chart.]

{\sl Remark}: When we apply the last bit-strings comparison for a more complex
holon tree, where there are ``repeating,'' ``empty'' and ``dispatching'' nodes,
these nodes also have Dewey labels. Of course, they must be discarded from the
computed substring set.

\section{The ``tendency to integration'' feature revisited}

At this point, it is worthwhile to return to the holon's tendency to
integration.

In the author's opinion, rejection of the structured programming principles is
often based upon efficiency of the running program. People can argue that
structured programs, even if they work correctly, will look like laboratory
prototypes where you can discern all the individual components, but which are
not daily usable. Building ``integrated'' products is an engineering principle
as valuable as structuring the design process.

In civil engineering design, it is presently a mandatory concept (that is known
as the ``Shanley Design Criterion'': collect several functions into one part).
We may see this principle working by comparing ``web'' iron structured bridges,
built in the eve of this century, to integrated pre-stressed concrete bridges.

Another example can be found in the rocket design: if you make a cross-section
of, for instance, the German V-2, you find, going from the outer to the center
of the section: external skin, structural rods, tank wall\dots; if you cut
across the Saturn-B moon rocket, you only find an external skin which is at the
same time a structural component and the tank wall. [Rocketry engineers have
used the ``Shanley Principle'' thoroughly when they use the fuel pressure
inside the tank to improve the rigidity of the external skin!]

In modern engineering, the implementation of a system is a modular one when
there is fear of ``wear and tear'' upon some parts. This kind of fear does not
exist for a computer program, what happens is the need of ``rewriting'' some
components because we understand the problem better or because the nature of
the program application changes.

With the holon concept, a holon can be modified; but then, we rebuild
(automatically) all the program in order to have a well-integrated program.

\bigskip

\noindent[In the appendix, we give another example of design integration.]

%************************************************
\chapter{User's View and System Implementation}
%************************************************

The holon system has two components:
\begin{enumerate}
    \item the analyzer, and
    \item the synthesizer.
\end{enumerate}

\section{The analyzer}

\noindent This component collects the holons written by the programmer and
builds the binary tree representation of the program. This binary tree is the
``logical'' support of the program and it is used to help the programmer to
design the program. For instance, it permits program interpretative execution.
The programmer can delete, append, change holons; his program binary tree is
kept in long-term storage as long as the program design is in progress.

The analyzer checks the syntactic correctness of the programmer's holons.
If there are several programmers, holons can be protected from illegal
alterations.

The analyzer outputs a binary tree in a standard format which is a standard one
for input to a synthesizer. The analyzer points out the recursive holons.

\section{The synthesizer}

\noindent The input of a synthesizer is a complete binary tree for a program,
the terminal holons of which are all written in terminal languages. It is
forbidden to synthesize a holon program where there are undefined holons.

The synthesizer produces a machine-code version of the program, using a code
generator for the used terminal languages and a ``synthesis strategy'' to
choose the implementation unit for each holon.

The synthesizer chooses the code generator and the synthesis strategy according
to information given by the programmer about the computer which will run the
program. (This is only a library retrieval program.)

The output of the synthesizer is a machine-code version of the program, ready
for execution on a computer.

\section{System control language}

\noindent The programmer uses a language composed of short command reserve
words to control the analyzer and the synthesizer.
A description of the syntax of such a language is given in a following chapter.

%************************************************
\chapter{The Initialization}
%************************************************

The Data Structure Command Syntax is an attempt to introduce levels in the data
structure as we have levels for the operations.

We shall first expound the different programming problems which have led to
this syntax design; next, we shall describe the syntax in a following section.

\section{Abstract data structures}

We want to allow the programmer to use ``Abstract'' data structures; i.e., at
some level in his program, he declares data structures belonging to an abstract
type.

An ``Abstract Type'' is a data structure type, the name of which is meaningful
in the programmer's mind; but, an abstract type doesn't automatically imply a
machine implementation.

Example: Suppose we are designing a chess playing program; and we use
extensively the concept of a ``chess-board pattern,'' i.e., a pattern of the
chess elements. We design a set of operations like: ``Evaluate Chess-Pattern,''
``Strengthen Chess-Pattern,'' etc\dots At this level in the program design, the
machine representation of this pattern is irrelevant but we need the concept of
an ``abstract'' data structure: Chess-Pattern. And according to the operation
which we shall perform upon a chess-pattern, we will design (refine) the
abstract data structure to a machine representation.
(The point is that the obvious implementation of a chess-board by an array
would probably be a mistake. For instance, it seems that chess-players memorize
the chess-board in the form of a list structure.\autocite{simon1969sciences}
Another interesting example is a checker program designed by Strachey; the
method used to select the checker-board representation is a revealing
one.\autocite{strachey1966system})

Creation and refining are allowed in the Holon Programming Language by commands
{\sl create} and {\sl refine}. (See further).

\section{Data structure declaration and scope}
\label{sec:data-structure}

We have adopted the declaration and scope used in {\sc algol}; nevertheless, in
top-down programming the {\sc algol} rules must be extended.

{\emergencystretch=1em When we are involved in the refinement of some operation
(level), the need of
new variables appears. But, generally the newly-introduced variables have a
scope larger than the scope of the level where the need of introductions
appears to the programmer.\par}

Thus, the design of some inner levels implies the alteration of the
declarations of an outer level. (An example of this programming behavior can be
noted in programming by stepwise refinements.) At the International Computing
Symposium '73 at Davos, Prof. Wirth used for displaying the methodology the
following device: Each refinement is written in one color on a transparency,
and by stacking the transparencies, we obtain the final program. But, what is
revealing in the final multi-colored program is that there are some colors
interspersing and not a genuine embedding of colors. I suppose that the program
will be reprinted in the Symposium proceedings. Another example may be seen in
Knuth's letter 1, p 7.\autocite{knuth1973review}

The problem with this kind of programming behavior is how do you manage
alterations of the program? Obviously, if you decide to delete a level, you
must delete all the level operations and carefully scan all the ``side''
modifications implied by this level.

In {\sc hpl},\marginnote{{\sc hpl}: Holon Programming Language} it is possible
to write declarations, the scope of which is not limited to the scope of the
holon where the declaration is written. The programmer declares a variable and
specifies the scope by naming an outer holon. The scope of this holon will be
the scope of the declared variable. This command is the \kw{declare at level}
command.

This command works only if the holon which contains the declaration statement
in its body belongs to the subtree of the holon which defines the variable
scope.

% Figure 10
\begin{marginfigure}[3cm]
\begin{tikzpicture}%
  [vertex/.style={circle,semithick,draw=black,fill=white},,
   node distance=2cm, minimum size=20pt,
   on grid,
   >=latex,>={Stealth[inset=0pt,length=6pt,angle'=28,round]}
  ]
  \node[vertex] (A) {A};
  \node[vertex,below left=1.25cm and 0.5cm of A] (B) {B};
  \node[vertex,right=1.5cm of B] (C){C};
  \node[vertex,right=1.5cm of C] (E){E};
  \node[vertex,below left=1.25cm and 0.5cm of C] (D) {D};
  \node[vertex,right=1.5cm of D] (F){F};
  \node[vertex,below left=1.25cm and 0.5cm of D] (G) {G};
  \node[vertex,below left=1.25cm and 0.5cm of F] (H) {H};
  \draw[->]
    (A) edge (B)
    (B) edge (C)
    (C) edge (D)
    (C) edge (E)
    (D) edge (F)
    (D) edge (G)
    (F) edge (H);
\end{tikzpicture}
\caption{A holon with different scopes.}
\label{fig:declare}
\end{marginfigure}

In Figure~\ref{fig:declare}, if we are using a \kw{declare at level} command
inside holon G, the only levels which can appear in this command are D, C, A.
To declare a variable at levels B, E, F, and H is not allowed because variables
declared at these levels can't appear in statements written at level G. The
analyzer checks the correct use of this command.

Now, let us tackle a new problem: suppose in Figure~\ref{fig:declare} that we
are designing the body of the holon H. Suppose that we need to use a variable
declared local to holon G (The holon G was designed before holon H. The meaning
is that a ``mistake'' was made during the writing of holon G: We did not
understand that the scope of a variable must be larger than the scope of holon
G. (It is not exactly a mistake but a lack of program understanding.)
We shall use inside holon H, the \kw{same as} command, i.e., to declare a
variable the same as another variable in another holon.

The analyzer will move a ``virtual'' declaration of the variable to a holon
level which encompasses the two involved holons. (In Figure~\ref{fig:declare},
this holon level would be level C.)

The declaration at level C is a ``virtual'' one because there is no
modification of the body of holon G. The alteration of the variable scope is
local to C; it appears only if H is in the G ``context.'' For instance, see
Figure~\ref{fig:virtual}; in this holon, there are two appearances of holon G.
The \kw{same as} command must in this case qualify the G involving holon G
(for instance, naming G as G \kw{in} D). But the ``virtual'' declaration at
holon level C has no effect upon the declarations of G \kw{in} M.

% Figure 11
\begin{marginfigure}
\begin{tikzpicture}%
  [vertex/.style={circle,semithick,draw=black,fill=white},,
   node distance=2cm, minimum size=20pt,
   on grid,
   >=latex,>={Stealth[inset=0pt,length=6pt,angle'=28,round]}
  ]
  \node[vertex] (A) {A};
  \node[vertex,below left=1.25cm and 0.5cm of A] (B) {B};
  \node[vertex,below left=1.25cm and 0.5cm of B] (M) {M};
  \node[vertex,below left=1.25cm and 0.5cm of M] (G') {G};
  \node[vertex,right=1.5cm of B] (C){C};
  \node[vertex,right=1.5cm of C] (E){E};
  \node[vertex,below left=1.25cm and 0.5cm of C] (D) {D};
  \node[vertex,right=1.5cm of D] (F){F};
  \node[vertex,below left=1.25cm and 0.5cm of D] (G) {G};
  \node[vertex,below left=1.25cm and 0.5cm of F] (H) {H};
  \draw[->]
    (A) edge (B)
    (B) edge (C)
    (C) edge (D)
    (C) edge (E)
    (D) edge (F)
    (D) edge (G)
    (F) edge (H)
    (B) edge (M)
    (M) edge (G');
\end{tikzpicture}
\caption{A holon with a ``virtual'' declaration.}
\label{fig:virtual}
\end{marginfigure}

\section{The initialization}

Watching people writing structured programs, you will notice that most of the
time, they start the writing of a level like:
\begin{itemize}
    \item[] Level XYZ:
    \begin{itemize}
        \item[] initialization;
        \item[] do something;
        \item[] do another thing:
        \item[] $\vdots$
    \end{itemize}
\end{itemize}
They will agree that they write the sublevel ``Initialization'' because it may
happen that it may be used. But really, when they lay down the level name upon
the paper, they don't know.

If you plan to read the content of this sublevel before the others, the reading
will be meaningless. As a matter of fact, this initialization level
conglomerates a lot of operations which must be there in order to respect the
sequential execution of the program; but why they exist must be discovered by
the reading of the other levels' bodies.

In our opinion, as the programmer makes choices during the design process,
parts of the ``initialization'' are implied by these choices. Thus, there
exists a one-to-one correspondence between each part of the initialization and
a design choice. Furthermore, in our opinion, the design choices which involve
initialization are mainly data structure choices. The result is that in {\sc
hpl}, after a declaration you can write a reserved word (\kw{initial})
followed by a holon name. The execution of this holon name is performed
according to the scope assigned to the variable. This is possible even with
\kw{declare at level}.

\section{Restricting access to data structures}
\label{sec:restricting}

We will allow the programmer to restrict access to some data structure types to
a limited set of holons. Variables of the restricted access type can be used
only by the holons, the names of which appear in the set.

It is also allowed to restrict access for a variable. If this variable belongs
to a restricted access type, the set of the holons which can use this variable
must be a subset of the holon list for the restricted access type.

Example:\\
\kw{create} 1 file \kw{restricted} read; write; purge \kw{end}\\
\kw{declare} top-secret: file \kw{restricted} write; purge \kw{end}

The checking of all the restrictions is made by the analyzer.

\section{Local refinement}
\label{sec:local-refinement}

The refinement of an abstract data type may be a local process. For instance,
we come back to the abstract concept of a ``chess-board pattern'': the way to
refine this concept in terms of terminal data structures depends upon the
functions we want applied to this abstract concept.

The programmer designing the holon which will choose the next movement will
maybe use a list representation of the chess-board; but, obviously the choice
will seem unacceptable for the programmer writing the holon which prints a
picture of the chess-pattern, this latter programmer will be very happy to use
a simple array. Thus, it seems to me to be useful to allow ``local to a holon''
refinements. In some holons the global refinement is switched to a local one.

(In a previous report, I had called that a ``dedicated refine.'')

\bigskip\noindent
{\sl Remark}. In the first implementation of {\sc hpl}, I don't plan to
implement ``Local Refinement,'' because I have not found an acceptable answer
to the check involved by the local refinement. In short, if locally you change
the way an abstract type is refined, that implies that variables, the scope of
which encompasses the two refinements, must be ``trasnlated'' from one
refinement to the other. You may assume that the programmer must provide the
translating holon at the right place, but it seems more sensible that the
analyzer requests these translating holons and decides itself where they must
go. Think of two programmers writing in parallel two subtrees of a holon where
they make local refinements; the analyzer would be able to output the correct
``interface holon'' specifications.

%************************************************
\chapter{Data Structure Commands Syntax}
%************************************************

We shall describe the syntax\marginnote{In order to shorten the text, we don't
use a full BNF description. Most of the syntactic elements will be described
using plain English.} of:
\begin{enumerate}
    \item Type descriptors;
    \item Variable names;
    \item The set of commands to build the data structure environment of a
program.
\end{enumerate}

\section{Type descriptors}
\label{sec:type-descriptors}

The type of a variable can be:
\begin{enumerate}
    \item either one of the terminal types, i.e., a type belonging to one of
the terminal language; or
    \item a structured type.
\end{enumerate}
A structured type is an agglomeration of terminal types structured like a tree.
In this tree, only the ``leaves'' are terminal types. The programmer can give
names to the nodes; the names correspond to the rooted subtrees and they define
new types. The several levels of the tree are pointed out by the use of
numbers. Each time we step down one level in the tree type, the level number of
the new level must be one unit greater than the level number of the previous
level.

\medskip

\noindent Example:\\
1 Human-Body\\
\begin{tabular}{ l l l l l }
&2 Head &3 Eyes  &4 Color&\kw{character-string} \\
&       &3 Hair  &4 Number&\kw{integer}\\
&2 Arms &3 Number&\kw{integer}\\
&       &3 Length&\kw{integer}\\
&2 Legs &3 Number&\kw{integer}\\
&       &3 Length&\kw{integer}\\
&2 Trunk&3 Length&\kw{integer}\\
&       &3 Type  &\kw{character};\\
\end{tabular}

\medskip\noindent
The names of the type can be composed of characters and the dash symbol (-).

It is possible to define an Array of Structures, i.e., in the defined type, a
subtree appears several times at the same level. In this case, only one subtree
is described associated with a subscript. The subscript is an integer or an
integer parameter. The subscript is used in declarations to define the number
of subtrees of this type; it is also used to select one subtree when we make a
reference to a variable.

{\sl Remark}. (1) It is said that the subscript is an integer or an integer
parameter (not an integer variable). When we design a new type, we point out
only the fact that some subcomponents of the tree can be repeated.

Example:\\
1 \quad Matrix\\
\-\hspace{1cm}2 \quad Row (\%R\%)\\
\-\hspace{2cm}3 \quad element-in-row (\%E\%) \kw{real};

In a declaration, values must be given to the parameter; but the given value
can be the name a pre-declared integer variable.

Example:
\begin{itemize}
    \item[] \kw{declare} A, B, C: Matrix($R=N$, $E=M$);
    \item[] \kw{declare} Z: Matrix($R=S$, $E=4$)
\end{itemize}
(2) The introduction of Arrays complicates a little bit the working of the
\kw{create} command (discussed later).

\section{Variable names (or identifiers)}

A variable name is a string of characters and dash symbols; normally space
between two words is forbidden. The length is not limited.

When the type of the variable name is a structured one, the programmer can also
access the subcomponents of the structured type. That is made by the use of the
reserved symbol \kw{of}. This symbol is used to link components from the
leaves to the root of the structured type. Nevertheless, there is no
requirement to give all the links, but the described linkage must be
unambiguous.

For instance:\\
After \kw{declare} Dumont, Einstein, Leon: Human\_Body; we can use the
following variable subcomponent references:
\begin{itemize}
    \item[] Number \kw{of} Arms \kw{of} Dumont
    \item[] Type \kw{of} Leon
    \item[] Color \kw{of} Hair \kw{of} Einstein
\end{itemize}
These three variable references are unambiguous. It must be noticed that only
the last word of a variable subcomponent reference is a variable name; the
others (except of \kw{of}s) are type descriptors. The first word stands for a
terminal type, and the analyzer can check if the forthcoming operation is a
legal one for this terminal type.

\section{Commands}

\medskip\noindent
The following commands will be described:\\
\-\hspace{1cm}\sym{create}, \sym{refine}, \sym{declare}.

\subsection{\sym{create}}

The \sym{create} command is used to create new types in the form of structured
types or of association of a name to a terminal type. An example of the latter
is:

\medskip

\-\hspace{2cm} \kw{create} color: \kw{character-string};

\medskip

\noindent Here ``color'' is not a variable, it is a data type. Normally some
restrictions upon the kind of character-string for ``color'' would be added to
the create command;\autocite{hoare1972chapter} nevertheless, at this time we
have not found a suitable syntax for this domain restriction (suitable syntax
for the programmer {\sl and} the analyzer.

In the former form (creation of structured types), it is not required that the
programmer writes the structured type thoroughly (from the root to the terminal
type leaves). The programmer can define abstract types.

Example: \kw{create} 1 chess\_board;

The number ``1'' indicates that it is a structured type, but it is an abstract
type.

The semantics of this command can be explained like this: when we start the
writing of a program, the analyzer is aware that the legal type set is the
``union'' of the terminal types of the terminal languages. Using the create
command, we can rename terminal types with more ``convenient'' type names for
our program, or we can really ``stretch'' the type set by ``creating''
structured types or abstract structured ones. An important fact is that only
the root of a structured type is added to the type set. If one of
the~\nameref{sec:type-descriptors} is preceded by the reserved word
\kw{create}, it becomes an example of structured type creation, but the type
set is
augmented by only one element: human\_body. (It will be shown that a
\kw{declare} command must associate a variable name to a type belonging to the
type
set.)

\subsection{\sym{refine}}

The \sym{refine} command is used to transform an abstract structure
step-by-step into a fully-defined structure. The command syntax is:

\medskip

\-\hspace{0.8cm} \kw{refine} \sym{list of abstract structures} \kw{by}
\sym{type descriptor};

\medskip\noindent
It is possible to apply the same refinement to several abstract structures;
that is the reason for the syntax component \sym{list of abstract structures}.

The programmer must write a level number `1' for each element appearing in the
list, whatever the actual level number of the abstract structure element. The
level numbers of the refining types (which follow the \kw{by}) are level
numbers relative to the level number `1'.

Example:

Suppose we have created the following abstract type:

\medskip

\-\hspace{0.8cm} \kw{create} 1 chess\_board;

\medskip\noindent
Now, we made a first refinement:
\medskip

\-\hspace{0.8cm} \kw{refine} 1 chess\_board \kw{by} 2 before\_the\_move \\
\-\hspace{5.35cm} 2 after\_the\_move;

\medskip\noindent
Further in the program, we perform new refinements:

\medskip

\-\hspace{0.8cm} \kw{refine} 1 before\_the\_move, 1 after\_the\_move \kw{by}\\
\-\hspace{3cm} 2 white\_positions \quad 2 black\_positions;

\medskip\noindent
At this time, the chess-board data type becomes:
\begin{itemize}
    \item[1] chess\_board
    \begin{itemize}
        \item[2] before\_the\_move
        \begin{itemize}
            \item[3] white\_positions
            \item[3] black\_positions
        \end{itemize}
        \item[2] after\_the\_move
        \begin{itemize}
            \item[3] white\_positions
            \item[3] black\_positions;
        \end{itemize}
    \end{itemize}
\end{itemize}
The relative level numbers can be used to add components at one level; for
instance, suppose we have the following command:

\medskip

\-\hspace{0.8cm} \kw{create} 1 chess\_board\\
\-\hspace{3cm} 2 before\_the\_move;

\medskip\noindent
Using the refine command, we write:

\medskip

\-\hspace{0.8cm} \kw{refine} 1 chess\_board \kw{by} 2 after\_the\_move;\\

\medskip\noindent
The following rules must be respected (to avoid rejection of the command by the
analyzer): [we write the syntax like ``\kw{refine} A \kw{by} B;'' in order to
shorten the rule descriptions.]
\begin{enumerate}
    \item Elements of the A-list must have appeared before in the program in a
\kw{create} command or in the B-list of a \kw{refine} command. The ``before''
must be understood as: in a Holon already ``swallowed'' by the analyzer.
    \item The refinement must not be a recursive one.
\end{enumerate}
A data structure is fully defined when all the refinements have led to terminal
types.

Refinement of a data type can have a scope; i.e., that locally to some
specified holons, a type descriptor refinement will be overlaid by a local
refinement. This local refinement can be restricted to some variables. The
syntax of this command looks like this:

\medskip

\-\hspace{0.8cm} \kw{Inside} \sym{holon names list}\\
\-\hspace{3cm} [\kw{for} \sym{variable names list}] \kw{refine} A \kw{by} B;

\medskip\noindent
Example:
\medskip

\-\hspace{0.8cm} \kw{Inside} print chess configuration\\
\-\hspace{2cm} \kw{refine} 1 before\_the\_move, 1 after\_the\_move \kw{by}
B;\\
\-\hspace{3cm} 2 \kw{character array} [1:8, 1:8];

\medskip\noindent
Outside the ``print chess configuration'' this refinement does not apply.

\medskip\noindent
{\sl Note}: In the first implementation of the holon language, we don't plan to
implement the~\nameref{sec:local-refinement}.

\subsection{\sym{declare}}

The semantics of a declaration are quite simple: It makes an association
between a variable name and a type descriptor. But, it also implies the
definition of a scope for this variable.

In conventional languages, the scope of a variable is determined by the
position of the declaration of this variable amid the program text. In {\sc
hpl}, the rule is: A variable declaration appears where the need for this
variable appears for the first time; and the declaration command is designed to
explain what is the scope of the variable.

The different forms of the declare command are:
\begin{itemize}
    \item[(a)] The ``local declare'' or ``simple declare.'' This command means
that the scope of the declared variable is similar to the scope of the holon
which contains the declare command.

    Example:

    \-\hspace{0.4cm} \kw{declare} Spassky, Fischer: Chess\_Board;

    [The type descriptor can be an abstract type; and this declaration can
appear in the program even if the refinement of the abstract type to a
structure of terminal types is not finished.]

    The syntax does not require that the declarations are put together in front
of a holon unit. They may appear anywhere between two holon language
statements; the analyzer will gather the declaration in due place.

    The type descriptor appearing in a ``declare'' statement must have appeared
in a ``create'' earlier in the program.
    \item[(b)] The ``at level declare.'' The logical reason for this command is
expounded in~\nameref{sec:data-structure}.

    An example of this command:

    \-\hspace{0.4cm} \kw{declare} Einstein: human\_body \kw{at}
process\_input level;

    [``process input'' is a holon name.]

    The semantics of this statement require that the holon, at the level of
which the declaration is made, belongs to the ``ancestors'' of the holon where
the command appears.
    \item[(c)] The ``same as declare.'' The logical reason for this command is
expounded in~\nameref{sec:data-structure}. The semantics are also explained in
that section.

    Example:

\-\hspace{0.4cm} \kw{declare} blackboard:\\
\-\hspace{1cm}draft \kw{same as} scratch \kw{in} trial first solution;

This example assumes that ``draft'' is a type descriptor, and that ``scratch''
is a variable which exists in the holon called ``trial first solution.'' The
type descriptor of the ``scratch'' variable must be the same as the one of the
declared variable; this is checked by the analyzer. [It is required that the
variable exists at the level of the holon, the name of which follows the
reserved word \kw{in}. It is not required that the variable is declared at
this holon level.]
\end{itemize}
Any of the three forms of declaration can contain an \kw{initial} command. The
following syntactic unit must appear before the connective symbol of the
declaration:
\medskip

\-\hspace{0.8cm} \kw{initial} \sym{holon name}

\medskip\noindent
Example:

\medskip

\-\hspace{0.8cm} \kw{declare} Fischer: chess\_board \kw{at} starting game
\kw{level initial}\\
\-\hspace{2cm} put chess set;

\subsection{The ``restricted'' option}

The ``restricted'' option is introduced before the connective symbol of either
a \kw{create} if we want to restrict access for data types or a \kw{declare}
if the restriction is for a variable.

If there is more than one holon in the restriction list, the end of the command
must be specified by an \kw{end}. See examples in~\nameref{sec:restricting}.

\section{The \texorpdfstring{\sym{value}}{<value>} option}

\medskip\noindent
It was said in the~\nameref{sec:boolean-holon} section that a boolean holon may
have access variables by ``value.'' In a boolean holon name, the variables
which we desire to be accessed by ``value'' are pointed out by writing the
reserved word \kw{value} before the identifier left bracket.

Example: Is matrix \%A\% with bounds \kw{value} \%N\% \kw{value} \%M\% a
square one?

This concludes the informal description of the ``Data Structure Commands.'' In
the next chapter, the data structure organization of a holon program will be
described.

%************************************************
\chapter{Data Organization in Holon Programming}
%************************************************

In this section, we shall explain how the data structures are introduced inside
the holon tree. The goal is to check if each holon appearance in a tree meets a
suitable data structure environment.

\section{Data structure ordering}

A holon has access to two kinds of data structures:
\begin{itemize}
    \item[(a)] Local data structures
    \item[(b)] Non-local data structures.
\end{itemize}
The elements belonging to these sets are issued from different kinds of
declarations.

We can order the local data structures from holon X:
\begin{enumerate}
    \item Declared in the holon X body;
    \item Declared from holons inner to X;
    \item Induced at the holon X level by \kw{same as} declarations.
\end{enumerate}

We can order the non-local data structures for holon X:
\begin{enumerate}
    \item Declared in an ``ancestor'' of X;
    \item Declared non-local from holon X (by a \kw{declare at level});
    \item Produced by a \kw{same as} declaration.
\end{enumerate}

The first holon of a program has only local data. A data structure
``non-local'' to a holon must appear ``local'' in an ancestor of this holon.

\section{The ``initial'' operation}

\medskip\noindent
Data structures, the declaration of which is followed by the ``initial''
option, generate links to the initial holons corresponding to these options.

For instance, the following holon text leads to the binary tree representation
in Figure~\ref{fig:initial}.

\medskip

\-\hspace{0.8cm} A\\
\-\hspace{0.8cm} \kw{begin} B; \kw{declare} M; C;\\
\-\hspace{2cm} \kw{declare} N \kw{initial} NI;\\
\-\hspace{2.5cm} D; \kw{declare} Q \kw{initial} QI;\\
\-\hspace{1.2cm} \kw{end}\\

\medskip
\noindent
[The test is a shortening of another one; B, C, D, NI, QI stand for holon
names; the type descriptors of the declared variables are omitted.]

% Figure 12
\begin{marginfigure}
\begin{tikzpicture}%
  [vertex/.style={circle,semithick,draw=black,fill=white},
  ell/.style = {% E as ellipse
    ellipse,semithick,draw=black,fill=white},
   node distance=2cm, minimum size=20pt,
   on grid,
   >=latex,>={Stealth[inset=0pt,length=6pt,angle'=28,round]}
  ]
  \node[vertex] (A) {A};
  \node[ell,below left=1.25cm and 0.5cm of A] (I) {INIT};
  \node[ell,right=2.5cm of A] (V) {\textit{var}: M,N,Q};
  \node[ell,below left=1.25cm and 0.5cm of I] (NI) {NI};
  \node[ell,right=1.6cm of NI] (QI){QI};
  \node[vertex,right=1.6cm of I] (B){B};
  \node[vertex,right=1.25cm of B] (C){C};
  \node[vertex,right=1.25cm of C] (D){D};

  \draw[->]
    (A) edge (I)
    (I) edge (B)
    (B) edge (C)
    (C) edge (D)
    (I) edge (NI)
    (NI) edge (QI);
  \draw[dashed,->] (A) -- (V);
\end{tikzpicture}
\caption{A holon with data structures declared after the ``initial''
operation.}
\label{fig:initial}
\end{marginfigure}

\section{Permanent and non-permanent parts of a holon}

We call the permanent part of a holon:
\begin{enumerate}
    \item The holon body,
    \item The local data structures set.
\end{enumerate}

We call the non-permanent part:
\begin{enumerate}
    \item The non-local data structures set.
\end{enumerate}

The permanent part has this name because whenever we write the name of a holon
inside a holon language statement, this permanent part can be attached to the
binary tree.

The non-permanent part describes the environment required by the holon; the
description of the environment is a list of variable names with a type
descriptor. The holon requires finding the same list (except variable names
which can be different) in any place where its name appears. The name
non-permanent comes from the fact that the variable names can change.

\medskip\noindent
{\sl Remark}. The checking of the consistency of the environment with the
permanent part is really a datatest operation; but this checking can be
performed completely by the analyzer (at ``compile-time'').

\section{Operations on a program structured by holons}

The operations that a programmer can do during program design are:
\begin{enumerate}
    \item To define a new holon;
    \item To insert in a holon statement an already existing holon which may
be:
    \begin{enumerate}
        \item Fully defined (all its pending nodes are written in terminal
languages);
        \item Partly defined;
    \end{enumerate}
    \item To write refinements (definitions of inner holons) of a holon, the
name of which appears in more than one place in the holon tree;
    \item To delete from the tree:
    \begin{enumerate}
        \item A holon name appearance in a holon statement;
        \item A holon definition (without ``erasing'' the holon name);
        \item A holon (name and definition).
    \end{enumerate}
\end{enumerate}

\section{Use of the non-permanent part of a holon}

When any of the previous operations occurs during the design of a program, the
analyzer uses the information gathered in the non-permanent part to check the
environment state required by the operation. The analyzer reports to the
programmer the reasons why the environment can preclude the completion of an
operation.

Each holon name node of the binary tree is linked to a non-permanent part
description which is used by the analyzer.

%************************************************
\chapter{Running the Synthesizer}
%************************************************

We shall expound by the use of a simple example how a synthesizer works.
Suppose the holon program, depicted in Figure~\ref{fig:analyzer}.
We use a triangle to point out a holon written in terminal language.

% Figure 13
\begin{figure}
    \centering
\begin{tikzpicture}%
  [vertex/.style={circle,semithick,draw=black,fill=white},
   triangle/.style={regular polygon,regular polygon
sides=3,semithick,draw=black,fill=white, minimum size=29pt},
   node distance=2cm, minimum size=20pt,
   on grid,
   >=latex,>={Stealth[inset=0pt,length=6pt,angle'=28,round]}
  ]
  \node[vertex] (A) {A};
  \node[vertex,below left=1.25cm and 0.5cm of A] (B) {B};
  \node[triangle,below left=1.25cm and 0.5cm of B] (E) {};
  \draw node [below left=1.25cm and 0.5cm of B] {E};
  \node[vertex,right=1.5cm of E] (F){F};
  \node[triangle,below left=1.25cm and 0.5cm of F] (G) {};
  \draw node[below left=1.25cm and 0.5cm of F] {G};
  \node[triangle,right=1.5cm of G] (H){};
  \draw node[right=1.5cm of G] {H};
  \node[vertex,right=3cm of B] (C){C};
  \node[triangle,right=1.5cm of C] (D){};
  \draw node[right=1.5cm of C] {D};
  \node[vertex,below left=1.25cm and 0.5cm of C] (B') {B};
  \node[vertex,right=1.5cm of B'] (I){I};
  \node[vertex,right=1.5cm of I] (B''){B};
  \node[vertex,below left=1.25cm and 0.5cm of I] (F'){F};
  \node[vertex,right=1.5cm of F'] (M){M};
  \node[triangle,below left=1.25cm and 0.5cm of M] (D'){};
  \draw node[below left=1.25cm and 0.5cm of M] {D};
  \node[triangle,right=1.5cm of D'] (E'){};
  \draw node[right=1.5cm of D'] {E};
  \draw[->]
    (A) edge (B)
    (B) edge (C)
    (B) edge (E)
    (E) edge (F)
    (F) edge (G)
    (G) edge (H)
    (C) edge (D)
    (C) edge (B')
    (B') edge (I)
    (I) edge (B'')
    (I) edge (F')
    (F') edge (M)
    (M) edge (D')
    (D') edge (E');
\end{tikzpicture}
    \caption{A holon program whose triangle-shaped nodes depict holons written
in a terminal language.}
    \label{fig:analyzer}
\end{figure}

The synthesizer generates the integrated program during two tree traversals.
These traversals will be described by a textual index pointing out the visited
nodes.

The first traversal is a postorder one. During this traversal, the synthesizer
generates intermediate code for the terminal holons. Moreover, it evaluates a
trade-off function, the parameters of which are the length of the intermediate
code, the length of an eventual procedure call for this holon, the frequency of
the holon in the program, and the kind of instructions surrounding the holon.
This evaluation decides the nature of the code integration: ``off-line'' code
(procedure) or ``in-line'' code (macro).

For instance, in Figure~\ref{fig:analyzer}, we will have the following sequence
of operations (the letters point out the visited nodes):
\begin{enumerate}
    \item E : Generate intermediate code;\\
              Trade-off function result: off-line (we suppose that).
    \item G : Generate intermediate code;\\
              trade-off: in-line.
    \item H : Generate intermediate code;\\
              trade-off: in-line.
    \item F : Trade-off: in-line (given F appears two times in the program,
this choice is possible if the sum of the code lengths of G and H is small).
    \item B : The intermediate code for B is composed of the calling sequence
of procedure B and the code for F. Given B appears three times, we suppose that
the trade-off result is: off-line.
    \item B : No-operation because the choice is already made.
    \item F : No-operation.
    \item D : Generate intermediate code;\\
              Trade-off: in-line.
    \item E : Generate a calling sequence for E.
    \item M : Trade-off: in-line.
    \item I : The trade-off gives immediately: in-line because ``I'' appears
only one time, thus the length of its code is irrelevant.
    \item B : No-operation.
    \item C : Trade-off: in-line. (C appears only one time).
    \item D : No-operation.
    \item A : We can start the second traversal.
\end{enumerate}
The second traversal is a preorder one. The synthesizer will generate the final
machine code using the information of the permanent and non-permanent data to
allocate storage and define the addresses or the references through a calling
sequence.

In the example, the first visited node (after the Node A) is node B. At this
point, we know that B is off-line, and we generate a calling sequence for B.
After that, we traverse the left subtree of B, building the code of the
procedure B. This code starts by a calling sequence of the procedure E,
followed by the code for this procedure. Next, the implementation of procedure
B resumes by the generation of the in-line machine code for F. Of course, code
to jump around the embedded procedures in generated.

The linear sequence of code generated for the holon A will look like:
$$c(B)\quad B:[\,c(E) \quad E: [\ ] GH\,]\quad c(B)\ GHD\ c(E)\ c(B)\ D$$

Meaning of the symbols:
\begin{itemize}
    \item[] $c(X)$ means call of procedure $X$;
    \item[] $X$: Points out the position of the code of procedure $X$, this
code is included between the two [ ];
    \item[] Y: Standing alone (not followed by a colon symbol), means an
in-line coding of the macro Y.
\end{itemize}
The two appearances of D in the symbol string don't mean that the generated
code at these places is similar; because the addresses used in the instructions
can be quite different.

Memory for the static data structures is allocated in a pseudo-stack with
overlapping of the data according to the scopes of the declarations.

The dynamic data structures (for instance, data involved in the recursive
procedures or data with dynamic bounds) are implemented by a stack.

{\emergencystretch=1em The synthesizer produces (during this second traversal)
the
machine-code-translation of the ``control switches'' used by the holon
statements. The behavior of the integrated program is driven by genuine machine
jumps and no more by the postorder traversal model used during the
analysis.\par}

\bigskip\noindent
{\sl Remark}. It must be noticed that the symbol string used to represent the
generated code contains only terminal holon names or off-line holon names. All
the other names used to indicate the program levels disappear.

%************************************************
\chapter{The Modified Syntax}
%************************************************

In this section, we shall describe some modifications to the syntax of the
holon language. These modifications don't fundamentally change the syntax; this
may be viewed as some alterations made to ease the handwriting of a holon
program. The relevancy of these changes depends upon the kind of equipment used
to submit programs to an analyzer. These modifications are for people using a
card input system. The content of this chapter is irrelevant if you use a
time-sharing system with a CRT display and some text editing program.

\section{The reserved word ``etc''}

The holon syntax requires that you write the holon name followed by the
definition (holon body) when you define a holon. Holon names can be very long
in order to be fully explanatory text.

The reserved word \kw{etc} permits one to (instead of writing a full holon
name) write only the beginning of the name followed by the word \kw{etc} and
the parameters if there are any.

For instance:
\begin{itemize}
    \item[] Invert matrix \%X\% with bounds \%N\% \%M\%
    \item[] \quad can be written
    \item[] Inv \kw{etc} \%X\% \%N\% \%M\%
\end{itemize}
There must be enough letters in the abbreviation to avoid ambiguities. In the
example, if the program contains several holon names beginning by ``Inv,'' the
submitted abbreviation is rejected by the analyzer. In the output document, the
input abbreviation is expanded.

\section{The text option}

The programmer can introduce a ``text'' as a part of a holon body. This must be
written:
\begin{center}
    \kw{text} \kw{begin} ``Any text ---'' \kw{end}
\end{center}
In this text, the programmer can explain the reason of some design decisions or
anything like this. In the holon body, the holon names embedded inside the
holon statements explain how the program works; in the ``text,'' the programmer
may expound why the program was written in this way. It is a ``comment'' to
help program readers to understand the logical basement of the program.

\section{Implicit holon names}

The holon syntax requires that, in statements as ``do while,'' ``repeat
until,'' ``case''---you write the name of one holon for the repeated part of
the statement.

Example:
\begin{itemize}
\setlength\itemsep{-0.3em}
    \item[] A
    \item[] \kw{begin} B; \kw{do} C \kw{while} Q; Z; \kw{end}
    \item[] C
    \item[] \kw{begin} CA; CB; CC; \kw{end}
\end{itemize}
This levelling can be omitted by writing explicitly the definition of the
``repeat'' part in the statement.

The previous example becomes:
\begin{itemize}
\setlength\itemsep{-0.3em}
    \item[] A
    \item[] \kw{begin} B;
    \item[] \quad\quad\quad \kw{do} \kw{begin} CA; CB; CC; \kw{end} \kw{while}
Q; Z; \kw{end}
\end{itemize}
The \kw{begin} and \kw{end} words are used by the analyzer to build correct
linkages in
the binary tree.

Holon names ``bracketed'' by \kw{begin} ---\ \kw{end} can appear between the
following
symbols:
\begin{itemize}
\setlength\itemsep{-0.2em}
    \item[] \kw{repeat} ---\ \kw{until}
    \item[] \kw{do} ---\ \kw{while}
    \item[] \kw{do} ---\ ; \qquad (if we have a statement \kw{while} ---\
\kw{do} ---\ ;)
    \item[] \kw{then} ---\ \kw{else}
    \item[] \kw{else} ---\ ;
    \item[] \kw{then} ---\ \kw{eor}
\end{itemize}
Only holon names can be used between \kw{begin} ---\ \kw{end}, not holon
statements.
Likewise, holon names for terminals can be omitted and the statement in
terminal language is ``bracketed'' by the \# --- and \#\# symbols.

\section{The ``proc'' and ``mac'' options}

The programmer can supersede the decision of the synthesizer by the use of the
\kw{proc} or \kw{mac} options.

A holon name followed by the reserved word \kw{proc} (\kw{mac}) means that
the programmer wants the implementation unit for this holon to be a procedure
(macro), regardless of the synthesizer's choice.

If this reserved word appears in the holon definition, the choice is valid for
all appearances of this holon in the program. If the reserved word follows a
holon name inside a holon statement, the choice is valid only at this place.

However, a \kw{mac} option for a recursive holon will always be rejected by
the synthesizer.

This option allows the programmer to ``contest'' the trade-off function of the
synthesizer. It may also be used during a new synthesizer design to test the
trade-off function.

%************************************************
\chapter{Holon System Control Language}
%************************************************

We shall describe a holon system control language. This language is an
``interface'' between the programmer and the implemented version of the
analyzer and synthesizer which he is using.

\section{Analyzer commands}

The majority of the control language commands deal with the analyzer.
\begin{itemize}
    \item[(A)] To start the design of a new program, the programmer uses the
following command:

    \smallskip
    \-\hspace{0.3cm}\kw{new program} \sym{holon name} \\
    \-\hspace{0.6cm}\kw{by} \sym{programmer identification} \kw{end}

    \smallskip
    The result of this command will be the opening of a file to store the
future holons of the program and the introduction of a new item in the file
which lists the undertaken programs.
    \item[(B)] To access a program, the programmer issues the \kw{program}
command:

    \smallskip
    \-\hspace{0.3cm}\kw{program} \sym{holon name} \\
    \-\hspace{0.6cm}\kw{by} \sym{programmer identification} \kw{end}

    \smallskip
    If the holon name is recorded as a ``program'' and if the programmer
identification is a legal one, the analyzer is ready to obey the programmer
requests.
    \item[(C)] The \kw{append} command is used to introduce the definitions of
new holons:

    \smallskip
    \-\hspace{0.3cm}\kw{append} \sym{holon} \dots \kw{end} \kw{append};

    \smallskip
    If holons in the ``append'' command have names which don't appear in holon
statements of other holons, the analyzer keeps these holons in a ``wait list''
for future use. Each time the analyzer finishes an append session for a
program, it tries to append the holons which are in the waiting list.

    After that, the programmer receives the following information:
    \begin{enumerate}
        \item An edited text of the appended holons;
        \item A report on syntactic errors;
        \item The status of the waiting list;
        \item A list of the undefined holons.
    \end{enumerate}
    \item[(D)] The programmer can used the \kw{undefine} command to erase the
last definition given for a holon. If the erased definition contains names of
already-defined holons, these latter definitions are not destroyed if these
holon names appear elsewhere in the program. The command is:

    \smallskip
    \-\hspace{0.3cm}\kw{undefine} \sym{holon name} \dots \kw{end}
\kw{undefine};

    \smallskip
    After an ``undefine'' session, the analyzer reports the list of the
undefined holons for the program.

    The appearance of the name of a holon existing in the wait list results in
the complete deletion of this holon: the definition and the name of this holon
are deleted.
    \item[(E)] The \kw{change} command is a combination of the two previous
ones:

    \smallskip
    \-\hspace{0.3cm}\kw{change} \sym{holon} \dots \kw{end} \kw{change};

    \smallskip
    For each holon appearing in the ``change'' list, the operations are the
following ones:
    \begin{enumerate}
        \item The previous definition of this holon becomes ``undefined.'' But
this operation is not conducted like the regular \kw{undefine} command: if the
erased definition contains names of already-defined holons, these latter
definitions are not destroyed if these holon names appear elsewhere in the
program or if they appear in the new definition.
        \item The new definition is appended.
    \end{enumerate}
    The use of this command is mainly the ``remodelling'' of the holon
statements of a holon.
    \item[(F)] The following command:

    \smallskip
    \-\hspace{0.3cm}\kw{output analyzed program} \sym{holon name} \kw{upon}\\
    \-\hspace{0.6cm}\sym{device description} \kw{by} \sym{programmer
identification};

    \smallskip
    is used to obtain the output (according to a standard format) of a program
(partly or completely finished). The device description points out the support
of the output: cards, tape, disk, etc.
    \item[(G)] The reverse operation of the previous one is made by:

    \smallskip
    \-\hspace{0.3cm}\kw{program} \sym{holon name} \kw{from} \sym{device
description}\\
    \-\hspace{0.6cm}\kw{by} \sym{programmer identification};
\end{itemize}

\section{Synthesizer commands}

\medskip\noindent
There are two commands:
\begin{itemize}
    \item[(A)] The first one is for requesting a program synthesis. The
programmer gives the name of the requested integration strategy:

    \smallskip
    \-\hspace{0.3cm}\kw{synthesize} \sym{program name} \kw{with}
\sym{strategy name} \\
    \-\hspace{0.6cm}\kw{by} \sym{programmer identification};

    \smallskip
    \item[(B)] The second one requests the execution of a synthesized program:

    \smallskip
    \-\hspace{0.3cm}\kw{execute} \sym{program name}\\
    \-\hspace{0.6cm}\kw{by} \sym{programmer identification};

    \smallskip
\end{itemize}
The \sym{programmer identification} is used in all these commands to protect
the program from illegal users.

%************************************************
\chapter{Program Example}
\label{ch:prog}
%************************************************

In this section, we shall show an example of a program written in the holon
programming language.

\section{Problem description}

This problem appeared in Dijkstra, Section 16.\autocite{Dijkstra1970-ja} We
quote here his problem description:
\begin{quotation}
    We consider a character set consisting of letters, a space and a point.
Words consists of one or more, but at most twenty letters. An input text
consists of one or more words, separated from each other by one or more spaces
and terminated by zero or more spaces followed by a point. With the character
valued function RNC (Read Next Character) the input text should be read from
and including the first letter of the first words up to and including the
terminating point. An output text has to be produced using the primitive
PNC($x$) (i.e., Print Next Character) with a character valued parameter.\dots
The text is to be subjected to the following transformation:
    \begin{enumerate}
        \item In the output text, successive words have to be separated by a
single space
        \item In the output text, the last word has to be followed by a single
point
        \item When we number the words 0, 1, 2, 3,\,\dots in the order from
left to right (i.e., in which they are scanned by repeated evaluation of RNC),
the words with an even ordinal number have to be copied, while the letters of
the words with an odd ordinal number have to be printed in the reverse order.
    \end{enumerate}
    For instance (using "\textvisiblespace" to represent a space) the input
text

    \smallskip
    \-\hspace{0.1cm}`this\textvisiblespace \textvisiblespace \textvisiblespace
is\textvisiblespace \textvisiblespace a\textvisiblespace silly\textvisiblespace
\textvisiblespace program\textvisiblespace \textvisiblespace.'

    \smallskip\noindent
    has to be transformed into

    \smallskip
    \-\hspace{0.1cm}`this\textvisiblespace si\textvisiblespace
a\textvisiblespace yllis\textvisiblespace program.'
\end{quotation}

\section{The holon program}

The following pages give the holon program written using the modified syntax.
The holons are written in the order of their design.\clearpage

\noindent odd inversion program\\
\-\hspace{0.2cm} \kw{begin} \\
\-\hspace{0.4cm} find first word starting character;\\
\-\hspace{0.4cm} \kw{repeat} find a word and print correctly \kw{until} end of
useful file;\\
\-\hspace{0.2cm} \kw{end}

\bigskip

\noindent find first word \kw{etc} \\
\-\hspace{0.2cm} \kw{begin} \\
\-\hspace{0.4cm} read first symbol;\\
\-\hspace{0.4cm} \kw{while} last read symbol is a space \kw{do} read next
symbol;\\
\-\hspace{0.2cm} \kw{end}

\bigskip

\noindent find a word \kw{etc} \\
\-\hspace{0.2cm} \kw{begin} \\
\-\hspace{0.4cm} buffer useful characters;\\
\-\hspace{0.4cm} print left to right or right to left according to even or
odd;\\
\-\hspace{0.4cm} find word terminator and print it and set end of file if full
stop;\\
\-\hspace{0.2cm} \kw{end}

\bigskip

\noindent buffer useful \kw{etc} \\
\-\hspace{0.2cm} \kw{begin} \\
\-\hspace{0.4cm} \kw{repeat} \kw{begin} buffer last read symbol;\\
\-\hspace{0.7cm} read next symbol;\\
\-\hspace{0.7cm} \kw{end} \\
\-\hspace{0.4cm} \kw{until} last read symbol is a space or a full stop;\\
\-\hspace{0.2cm} \kw{end}

\bigskip

\noindent find word terminator \kw{etc} \\
\-\hspace{0.2cm} \kw{begin} \\
\-\hspace{0.4cm} \kw{if} last read symbol is a space \kw{then} skip other
spaces;\\
\-\hspace{0.4cm} \kw{if} last read symbol is a full stop\\
\-\hspace{0.7cm} \kw{then} print it and set end of file;\\
\-\hspace{0.7cm} \kw{else} print a space;\\
\-\hspace{0.2cm} \kw{end}

\bigskip

\noindent print left \kw{etc} \\
\-\hspace{0.2cm} \kw{begin} \\
\-\hspace{0.4cm} \kw{if} even word\\
\-\hspace{0.7cm} \kw{then} print left-right;\\
\-\hspace{0.7cm} \kw{else} print right-left;\\
\-\hspace{0.4cm} increase word count;\\
\-\hspace{0.2cm} \kw{end}

\clearpage

\noindent even word\\
\-\hspace{0.2cm} \kw{begin} \\
\-\hspace{0.4cm} \kw{declare} even: \kw{boolean at} odd inversion program
\kw{level}\\
\-\hspace{0.4cm} \kw{initial} \# even $\gets$ {\it true} \#\#;\\
\-\hspace{0.4cm} \# even $\gets$ {\it true} \#\#;\\
\-\hspace{0.2cm} \kw{end}

\bigskip

\noindent increase word count\\
\-\hspace{0.2cm} \kw{begin} \\
\-\hspace{0.4cm} \# even $\gets$ \kw{not} even \#\#;\\
\-\hspace{0.2cm} \kw{end}

\bigskip

\noindent skip other spaces\\
\-\hspace{0.2cm} \kw{begin} \\
\-\hspace{0.4cm} \kw{repeat} read next symbol \kw{until} last read symbol is
not a space;\\
\-\hspace{0.2cm} \kw{end}

\bigskip

\noindent read first symbol\\
\-\hspace{0.2cm} \kw{begin} \\
\-\hspace{0.4cm} \kw{declare} lrs: \kw{character at} odd inversion program
\kw{level};\\
\-\hspace{0.4cm} \# lrs $\gets$ RNC \#\#;\\
\-\hspace{0.2cm} \kw{end}

\bigskip

\noindent last read symbol is a space\\
\-\hspace{0.2cm} \# lrs $\equiv$ '\textvisiblespace\,' \#\#;

\bigskip

\noindent read next symbol\\
\-\hspace{0.2cm} \# lrs $\gets$ RNC \#\#;

\bigskip

\noindent end of useful file\\
\-\hspace{0.2cm} \kw{begin} \\
\-\hspace{0.4cm} \kw{declare} EOF: \kw{boolean at} odd \kw{etc level};\\
\-\hspace{0.4cm} \kw{initial} \# EOF $\gets$ {\it false} \#\#;\\
\-\hspace{0.4cm} \# EOF $\equiv$ {\it true} \#\#;\\
\-\hspace{0.2cm} \kw{end}

\bigskip

\noindent buffer last read symbol\\
\-\hspace{0.2cm} \kw{datatest} number of symbols in the buffer less than 20;
\kw{end}\\
\-\hspace{0.2cm} \kw{begin} \\
\-\hspace{0.4cm} \kw{declare} buffer: \kw{character} (20) \kw{at} find a
word \kw{etc level};\\
\-\hspace{0.4cm} \kw{declare} pointer: \kw{integer at} find a word \kw{etc
level};\\
\-\hspace{0.4cm} \kw{initial} \# pointer $\gets 0$ \#\#;\\
\-\hspace{0.4cm} \# pointer $\gets$ pointer $+$ $1$; buffer(pointer) $\gets$
lrs \#\#;\\
\-\hspace{0.2cm} \kw{end}

\clearpage

\noindent last read symbol is a space or a full stop\\
\-\hspace{0.2cm} \# (lrs $\equiv$ '\textvisiblespace\,') \kw{or} (lrs $\equiv$
'.\,') \#\#

\bigskip

\noindent last read symbol is not a space\\
\-\hspace{0.2cm} \# lrs $\neq$ '\textvisiblespace\,' \#\#

\bigskip

\noindent last read symbol is a full stop\\
\-\hspace{0.2cm} \# lrs $\equiv$ '.\,' \#\#

\bigskip

\noindent print it and set \kw{etc} \\
\-\hspace{0.2cm} \# PNC('.\,'); EOF $\gets$ {\it true} \#\#

\bigskip

\noindent print a space\\
\-\hspace{0.2cm} \# PNC('\textvisiblespace\,') \#\#

\bigskip

\noindent print left-right\\
\-\hspace{0.2cm} \kw{begin} \\
\-\hspace{0.4cm} \kw{declare} point2: \kw{integer};\\
\-\hspace{0.4cm} \kw{initial} \# point2 $\gets$ $1$ \#\#;\\
\-\hspace{0.4cm} \kw{repeat} \# PNC(buffer(point)); point2 $\gets$ point2 $+$
$1$
\#\#\\
\-\hspace{0.4cm} \kw{until} \# point2 $>$ pointer \#\#;\\
\-\hspace{0.2cm} \kw{end}

\bigskip
\noindent print right-left\\
\-\hspace{0.2cm} \kw{begin} \\
\-\hspace{0.4cm} \kw{while} \# pointer $\neq$ $0$ \#\# \kw{do} \\
\-\hspace{0.7cm} \kw{begin} \\
\-\hspace{1cm}  \# PNC(buffer(pointer)); pointer $\gets$ pointer $-$ $1$
\#\#;\\
\-\hspace{0.7cm} \kw{end} \\
\-\hspace{0.2cm} \kw{end}

\section{The integrated program}

The synthesizer will produce code for a program equivalent to the one written
below. In this code, the block structure will be disentangled and the storage
for all the variables will be allocated on the pseudo-stack.

\bigskip\noindent
[This intermediate ALGOL-program form doesn't exist in the system; it is only
said that the generated code will be equivalent to the code generated for the
following program by an ALGOL compiler.]

\clearpage

\noindent\kw{begin} \\
\ind{1}\kw{declare} even: \kw{boolean}, lrs: \kw{character}, EOF: \kw{boolean};
\\
\ind{1}even $\gets$ {\sl true}; EOF $\gets$ {\sl false};\\
\ind{1}lrs $\gets$ RNC; \kw{while} lrs $\equiv$ \spchar\ \kw{do} lrs $\gets$
RNC;\\
\ind{1}\kw{repeat} \\
\ind{2}\kw{begin} \\
\ind{3}\kw{declare} buffer: \kw{character}($20$), pointer: \kw{integer};\\
\ind{3}pointer $\gets$ $0$;\\
\ind{3}\kw{repeat} \\
\ind{4}\kw{if} pointer $\geq$ $20$ \kw{then} (datatest exit);\\
\ind{4}pointer $\gets$ pointer $+$ $1$;\\
\ind{4}buffer(pointer) $\gets$ lrs;\\
\ind{4}lrs $\gets$ RNC;\\
\ind{3}\kw{until} (lrs $\equiv$ \spchar) \kw{or} (lrs $\equiv$ '.');\\
\ind{3}\kw{if} even $\equiv$ {\sl true}\\
\ind{4}\kw{then} \\
\ind{5}\kw{begin} \kw{declare} point2: \kw{integer}; point2 $\gets$ $1$;\\
\ind{5}\kw{repeat} PNC(buffer(point2)); point2 $\gets$ point2 $+$ $1$;\\
\ind{5}\kw{until} point2 $>$ pointer;\\
\ind{5}\kw{end} \\
\ind{4}\kw{else} \\
\ind{5}\kw{begin} \kw{while} pointer $\neq$ $0$ \kw{do} \\
\ind{6}\kw{begin} PNC(buffer(pointer)); pointer $\gets$ pointer $-$ $1$;
\kw{end} \\
\ind{5}\kw{end} \\
\ind{3}even $\gets$ \kw{not} even;\\
\ind{3}\kw{if} lrs $\equiv$ \spchar\\
\ind{4}\kw{then} \kw{begin} \kw{repeat} lrs $\gets$ RNC \kw{until} lrs $\neq$
\spchar; \kw{end} \\
\ind{3}\kw{if} lrs $\equiv$ '.'\\
\ind{4}\kw{then} \kw{begin} PNC('.'); EOF $\gets$ {\sl true}; \kw{end} \\
\ind{4}\kw{else} PNC(\spchar);\\
\ind{2}\kw{end} \\
\ind{1}\kw{until} EOF $\equiv$ {\sl true};\\
\noindent\kw{end}

\clearpage

\section{Comments}

A. I confess that I made a change after reading Dijkstra's students' answer:
the holons ``even word'' and ``increase word count'' were written:

\bigskip

\noindent even word\\
\-\hspace{0.2cm} \kw{begin} \\
\-\hspace{0.4cm} \kw{declare} wordcount: \kw{integer at} odd inversion
program \kw{level}\\
\-\hspace{0.4cm} \kw{initial} \# wordcount $\gets$ $0$ \#\#;\\
\-\hspace{0.4cm} \# even(wordcount) $\gets$ {\it true} \#\#;\\
\-\hspace{0.2cm} \kw{end}

\bigskip

\noindent increase word count\\
\-\hspace{0.2cm} \kw{begin} \# wordcount $\gets$ wordcount $+$ $1$ \#\#
\kw{end} \\

\medskip
\noindent
I had some trouble after writing this because I was afraid of an overflow of
``wordcount.'' That would require the introduction of a ``datatest'' with error
recovery in ``increase word count.'' The error recovery would be to ``reset''
the counter. But with this datatest, the holon ``increase wordcount'' seems to
be a very cumbersome one. It is not the overflow that really bothers me, but
that there may be an ``overflow interrupt.'' So I solve all of these problems
by writing this (retrospectively) incredible holon:

\bigskip
\noindent
increase word count\\
\-\hspace{0.2cm} \kw{begin} \# wordcount $\gets$ [wordcount $+$ $1$] \kw{mod}
$2$
\#\# \kw{end}

\bigskip
\noindent
It was the wordcount \kw{integer}\ declaration that hindered me from thinking
about a
boolean variable.

\medskip
\noindent
B. The reader may think that the result of the work is a very verbose program.
That is true, but the holon program explains the process in a different way
from a ``normal'' program:
\begin{enumerate}
    \item The logic of the program answer to the problem is contained in the
program text. You don't need to disentangle this logic from the instruction
sequence.
    \item The program text is self-explanatory.
\end{enumerate}

\medskip\noindent
C. Normally, with a real holon system, the programmer would submit his
first-written holons and receive (for instance, on punch cards which he shall
put in front of the definition) the list of the undefined holons.

\medskip\noindent
D. Advice on the usefulness of the system must not be established by the fast
that for so small a problem, the generation of the holon answer seems longer
than with a normal language. It is better to ask oneself: ``\dots and when
faced with a program a thousand times as large, should I compose it in the same
way?''\footnote{Ibid.}

\medskip\noindent
E. The use of the ``datatest'' must be noticed. It is put at its place as a
``fire-wall'' for the program.

Reading Knuth's version of the program suggests that it would be more sensible
to write a ``datatest'' with the error-recovery
option.\autocite{knuth1973review} The processing of the string will continue
even after a ``too long word.'' For instance:

\bigskip\noindent
\ind{1}\kw{datatest}\\
\ind{2}number of symbols in the buffer less than 20\\
\ind{2}\kw{else} \kw{begin} \\
\ind{3}print a quoted buffer;\\
\ind{3}reset pointer to zero;\\
\ind{2}\kw{end} \\
\ind{1}\kw{end}

\bigskip\noindent
With this datatest the following odd word:

\medskip
\ind{1}thiswordislongerthantwentycharacters

\medskip\noindent
will print as:

\medskip
\ind{1}``thiswordislongerthan''sretcarahcytnewt

\bigskip\noindent
[I forgot in the program to write the refinement of the datatest. Normally,
this must be written.]

%************************************************
\chapter[Macro-processing of Holon Name Parameters]{Macro-processing of Holon
\\ Name Parameters}
%************************************************

The previous-described syntax allows only the use of parameters for variables.
Parameters can be used for holon names provided we macro-process holon
definitions if there is a holon name parameter.

For instance, the postorder traversal of a binary tree is an algorithm which
appears frequently in this survey; the semantics of ``visit the node'' is the
only change. One may imagine a super-holon: ``traverse binary tree in postorder
using \%holon\% for visiting the node.''

Each appearance of this holon will produce an expansion of its body, in order
to build the whole program tree structure.

For the time being, we don't plan to implement this feature in the first
version of the holon language.

%************************************************
\chapter{Holon Parallel Programming}
%************************************************

Reading the book of Brinch Hansen\autocite{hansen1973operating} and Hoare's
paper,\autocite{hoare1972towards} it appears to me that it would be interesting
to study if, with some modifications, the holon language can be used to program
parallel processes. As a matter of fact, the language was not designed with
this purpose in mind.

I have not studied thoroughly this extension, at the present time.
Nevertheless, some of Hoare's requirements seem to be fulfilled more or less
easily:
\begin{enumerate}
    \item The \kw{cobegin} and \kw{coend} notation for disjoint processes.
Moreover, the non-permanent parts of the co-processed holons can be used to
check the disjointness at analyze-time.
    \item The \kw{exit} statement as it is defined (\kw{goto} the end of a
holon) fulfills the ban of jump out of a parallel statement (the parallel
statement is a co-holon).
    \item The ``when B'' condition would be a datatest with the error recovery
option: \kw{datatest} B \kw{else} waiting \kw{end}.
    \item The ``array remapping'' may be interpreted as a local refinement.
\end{enumerate}
This section is only a hint to an eventual extension of the holon language.

%************************************************
\chapter{Holon Operating System}
%************************************************

\section{Introduction}

We want to use the fact that we know the structure of a program in order to
manage a virtual memory system. Since we have a structured program, we have a
better knowledge of the behavior of the program. We can determine exactly the
future behavior of a program; we want to use this information to build a system
furnishing its working set beforehand to a program. Among the instructions of a
program, there will be some page transfer requests; but, those requests are
written automatically in due places in order to avoid page faults; the Holon
Operating System will prepare a suitable working set for each executed program.

\section{Paging of a holon program}

The terminal levels, linked by some suitable code corresponding to the machine
code for the holon language statements
(\kw{if} \dots\kw{then} \dots\kw{else} \dots\kw{do} \dots\kw{while}
\dots\kw{etc} \dots) are ``paged.'' These
pages correspond to the operational (or functional) part of the program; now,
we must also ``page'' the data structure, for that purpose we manage two
different memory zones: a paged stack and a paged pseudo-stack.

In the pseudo-stack, we will put all the variables---the dimensions and the
number of which may be known at compile time (this includes all the variables
except the ones appearing in arrays with dynamic bounds and in recursive
procedures). The variables in the pseudo-stack are ``overlaid'' if it is
possible; for instance the two component variables of an if-then-else clause
are overlaid. But the maximum length of this pseudo-stack may be known at
compile time. And at any point in the program, we may specify what is the
length of the used part of the pseudo-stack.

Now, in the stack we put all the data, the number or the dimensions of which
depend on the computation (dynamic arrays, data local to recursive procedures).

The difference between the two stacks is: we take two computation points in the
program (two different positions of a textual index), for instance $p_1$ and
$p_2$, with the following condition: $p_2$ is to the right of $p_1$ in the
holon tree (the meaning of which is: there are some computations to do for
moving the textual index from point $p_1$ to $p_2$); now we call $\psi(p_1)$
and $\psi(p_2)$ the extents of the pseudo-stack at point $p_1$ and $p_2$, and
in the same way S1 ($p_1$) and S1($p_2$); $\psi(\pi)$ is a fixed value not
correlated to the input data string; but S1($\pi$) depends on the input string.

Now, look at Figure~\ref{fig:op}---it represents a holon program with its code,
and the two ``paged'' stacks. (Remark: we suppose temporarily that there exists
only in-line code.)

% Figure 14
\begin{figure}
    \centering

\begin{tikzpicture}%
  [vertex/.style={circle,semithick,draw=black,fill=white},
   node distance=2cm, minimum size=20pt,
   on grid,
   >=latex,>={Stealth[inset=0pt,length=6pt,angle'=28,round]}
  ]
  \node[vertex] (A) {A};
  \node[vertex,below=1.2cm of A] (B) {B};
  \node[vertex,below=1.2cm of B] (E) {E};
  \node[vertex,right=1.2cm of E] (F) {F};
  \node[vertex,below=1.2cm of E] (one) {1};
  \node[vertex,below=1.2cm of F] (two) {2};
  \node[vertex,right=2.4cm of B] (C) {C};
  \node[vertex,below=1.2cm of C] (G) {G};
  \node[vertex,right=1.2cm of G] (H) {H};
  \node[vertex,below=1.2cm of G] (three) {3};
  \node[vertex,below=1.2cm of H] (four) {4};
  \node[vertex,right=2.4cm of C] (D) {D};
  \node[vertex,below=1.2cm of D] (I) {I};
  \node[vertex,right=1.2cm of I] (J) {J};
  \node[vertex,below=1.2cm of I] (five) {5};
  \node[vertex,below=1.2cm of J] (six) {6};
  \node[vertex,right=1.2cm of J] (K) {K};
  \node[vertex,below=1.2cm of K] (seven) {7};
  \draw[->]
    (A) edge (B)
    (B) edge (E)
    (E) edge (one)
    (E) edge (F)
    (F) edge (two)
    (B) edge (C)
    (C) edge (G)
    (G) edge (H)
    (G) edge (three)
    (H) edge (four)
    (C) edge (D)
    (D) edge (I)
    (I) edge (five)
    (I) edge (J)
    (J) edge (six)
    (J) edge (K)
    (K) edge (seven);

    % Code pages
    \node[below left=1cm and 0.4cm of one] (cptopleft) {};
    \node[right=8.4cm of cptopleft] (cptopright) {};
    \draw[semithick] (cptopleft.center)--(cptopright.center);

    \node[below=0.6cm of cptopleft] (cpbotleft) {};
    \node[right=8.4cm of cpbotleft] (cpbotright) {};
    \draw[semithick] (cpbotleft.center)--(cpbotright.center);

    \foreach \x in {0,...,14}
        \draw[semithick]
([xshift=\x*0.6cm]cptopleft.center)--([xshift=\x*0.6cm]cpbotleft.center);

    \node[below right=0.3cm and 0.6cm of
cptopright,align=left,font=\footnotesize\sffamily]
    {Code\\pages};
    % PS Pseudo-Stack Pages
    \node[below=0.55cm of cpbotleft] (pstopleft) {};
    \node[right=6.6cm of pstopleft] (pstopright) {};
    \draw[semithick] (pstopleft.center)--(pstopright.center);

    \node[below=0.6cm of pstopleft] (psbotleft) {};
    \node[right=6.6cm of psbotleft] (psbotright) {};
    \draw[semithick] (psbotleft.center)--(psbotright.center);

    \foreach \x in {0,...,11}
        \draw[semithick]
([xshift=\x*0.6cm]pstopleft.center)--([xshift=\x*0.6cm]psbotleft.center);

    \node[below right=0.3cm and 1.3cm of
pstopright,align=left,font=\footnotesize\sffamily]
    {PS Pseudo-stack\\pages};

    % S Stack Pages
    \node[below=0.55cm of psbotleft] (stopleft) {};
    \node[right=4.2cm of stopleft] (stopright) {};
    \draw[semithick] (stopleft.center)--(stopright.center);
    \node[below=0.6cm of stopleft] (sbotleft) {};
    \node[right=4.2cm of sbotleft] (sbotright) {};
    \draw[semithick] (sbotleft.center)--(sbotright.center);
    \foreach \x in {0,...,7}
        \draw[semithick]
([xshift=\x*0.6cm]stopleft.center)--([xshift=\x*0.6cm]sbotleft.center);
    % dashed extension
    \node[right=6.6cm of stopleft] (stopfarright) {};
    \node[right=6.6cm of sbotleft] (sbotfarright) {};
    \draw[semithick,dashed] (stopright.center)--(stopfarright.center);
    \draw[semithick,dashed] (sbotright.center)--(sbotfarright.center);

    \node[below right=0.3cm and 1.1cm of
stopfarright,align=left,font=\footnotesize\sffamily]
    {S Stack pages};
\end{tikzpicture}
    \caption{A holon program with its corresponding pages and ``paged''
stacks.}
    \label{fig:op}
\end{figure}

The pagination of the seven terminal holons is made like this: the code of the
terminal holons composes a string of machine code and we break the string,
putting the first $x$ words in the first page, the next $x$ words in the second
page, etc.; i.e., the paging isn't related to the structure of the program.

Now, in the next section, we shall see how a Holon Operating System would work.

\section{Monitoring the page transfers}

We shall call the ``working set for a holon'' the code and data pages needed by
a holon to avoid any disruption during the computation. The important feature
of the holon tree is that the ``working set'' is almost---because of
S---well-defined. For instance, in Figure~\ref{fig:op}, the working set for
holon A is composed of all the code and PS pages, plus an amount of S pages
(this amount is not known, but the holon's need of some space (some pages) in
the S stack may be known). Now, the working set of the holon B is composed of
the pages needed to perform the computation involved in the holon B, but also
the first pages needed by holon C to start its computation: we need the last
ones because we must avoid disruption not only during the holon computation,
but also between two holons.

As long as a holon computation progresses, the holon may also indicate the
pages it doesn't need anymore.

With these two operations in mind (a holon may make page requests in advance; a
holon may also indicate the pages to be dropped), it may be understood that we
may build a holon operating system where, as much as possible, the program's
behavior is managing the page requests.

\section{Structure of the computer}

We need a computer with an internal paged store, and (probably) with a
two-level back store. [The matter of the number of levels will appear further
on, because in the Holon Operating System it will appear that the page size
will be determined by the transfer speed between the levels of the store
hierarchy.]

We need also a ``monitor,'' this will be either a program or a small hardware
computer. The purpose of the monitor is: to receive the page requests from the
running programs and to manage the transfer of the pages.

\section{Fundamental structure of the paging algorithm}

[The following description is a sketch---it is the departure point of a study.]

We first suppose that the processor is allocated to the different processes
according to a scheduling policy. Each holon will request its working set
before starting execution. Nevertheless, the monitor will discard requests from
``overly ambitious'' holons. The definition of an ``overly ambitious'' holon is
up to a computer center manager---it will depend on the average job profile;
and, for instance, the ``overly ambitious'' barrier may be changed during the
day.

If a holon is discarded (because it is overambitious), the hierarchic levels of
this holon are stepped down until the first ``moderate'' holon is found; its
page requests are obeyed as soon as possible and this holon starts running.
Now, during its execution, it will request pages in advance and the monitor
will try to bring these pages in time; it is also the holon program that points
out the pages which may be overlaid by others, because it doesn't need them
anymore. It is only in case of emergency (deadlock) that the monitor decides to
remove a page not released by a holon.

\section{Working set for the fundamental instructions}

\subsection{If-then-else statement}

Suppose the statement `\kw{if} Q \kw{then} T \kw{else} F; next;' (`next' is not
necessarily
a holon name; it may be another holon statement).

Suppose that we have the following paginations: (We consider only ``code
pages'')
\begin{itemize}
    \item[] Q\,: $Q_1\ Q_2\ Q_3\cdots Q_N$ pages
    \item[] T\,: $T_1\ T_2\ T_3\cdots T_N$ pages
    \item[] F\,: $F_1\ F_2\ F_3\cdots F_N$ pages
    \item[] next\,: $N_1\ N_2\cdots N_X$ pages
\end{itemize}
The working set for the if-then-else statement is: ($Q_1\ Q_2\cdots T_1\ F_1$);
like this, the statement will be executed without disruption. The test is made
and the first page of both parts of the alternative are included; in the
beginning of the $T_1$ page, for instance, there will be requests to the
monitor:
\begin{enumerate}
    \item Page from $F_1$ may be overlaid,
    \item Please bring the following pages into the internal core:
$T_2,\allowbreak T_3,\allowbreak T_4,\ldots\allowbreak T_M,\allowbreak N_1$.
\end{enumerate}
Without waiting for the answer to the request, we start the execution of the
$T_1$ page instructions.

\medskip\noindent
{\sl Remark}. The requested pages depend on the nature of the holon T; if this
one is a new if-then-else statement, then the requested pages will correspond
to the if-then-else working set. In this case, we have brought into memory and
unused page (in the example: $F_1$), but that is not a problem, as we don't
want to reduce its minimum number of page transfers during a program
execution---we want to limit to a minimum the number of page faults.

\subsection{Case statement}

It appears that the case statement may be managed in the same way as the
if-then-else statement.

\subsection{Repeat and while statements}

Suppose statement `\kw{repeat} P \kw{until} Q; next;' with the pagination:
\begin{itemize}
    \item[] P\,: $P_1\ P_2\ P_3\cdots P_N$
    \item[] Q\,: $Q_1\ Q_2\ Q_3\cdots Q_N$
    \item[] next\,: $N_1\ N_2\ N_3\cdots N_X$
\end{itemize}
The working set is: ($P_1\ P_2\cdots P_N\ Q_1\cdots Q_N\ N_1$). In fact the
selected $P_i$ will depend on the $P$ statement type; but normally (except
tying up the internal memory) the deletion of pages will be delayed until
reaching the beginning of the page $N_1$.

Example: In the if-then-else statement, when we reach $Q_2$, we may request
deletion of $Q_1$; but if this statement is embedded in a repeat statement,
this deletion is not sensible, because it will overload the page transfer
mechanism. In the same way, embedded \kw{repeat} statements lead to fixed
pages.

\section{Maximum number of pages fixed}

In the previous section we presented a memory policy management that is
program-oriented, because it is the program itself that requests pages and also
deletes them according to its needs (the monitor preempts pages only if the
pages fixed by the processes tie up the internal memory.) There exists another
policy: to fix a maximum for the number of code pages allocated to a process.
Because we have structured programs, in front of each page we may, at
compile-time, write the list of the ($n-1$) pages requested to work ($n$ is the
maximum number allowed) and the list of page modifications for the pseudo-stack
and the stack. With this policy, we have a deletion each time we leave a page,
followed by a page transfer (except in the case where we are in a \kw{repeat}
statement, the extent of which is less than $n$; see the following example.

\begin{marginfigure}
\begin{tikzpicture}%
  [node distance=2cm, minimum size=20pt,
   on grid,
  ]
    \node (pstopleft) {};
    \node[right=6cm of pstopleft] (pstopright) {};
    \draw[semithick] (pstopleft.west)--(pstopright.west);

    \node[below=0.6cm of pstopleft] (psbotleft) {};
    \node[right=6cm of psbotleft] (psbotright) {};
    \draw[semithick] (psbotleft.west)--(psbotright.west);

    \foreach \x in {0,...,10}
        \draw[semithick]
([xshift=\x*0.6cm]pstopleft.west)--([xshift=\x*0.6cm]psbotleft.west);
    \foreach \x in {1,...,10}
        \node[xshift=\x*0.6cm-0.67cm,yshift=-.29cm] {\x};
\end{tikzpicture}
\caption{A program with 10 pages.}
\label{fig:pages}
\end{marginfigure}

In Figure~\ref{fig:pages}, pages $1$, $2$, and $3$ compose a \kw{repeat}
statement. We
list the page requests for $n=5$ on entry in:
\begin{itemize}
\setlength\itemsep{0.1em}
    \item[$1$:] $1,2,3,4,5$
    \item[$2$:] $2,3,1,4,5$
    \item[$3$:] $3,1,2,4,5$
    \item[$4$:] $4,5,6,7,8$
    \item[$5$:] $5,6,7,8,9$
\end{itemize}

This example shows that there is no transfer and deletion until we reach the
head of page 4: but at that time, we have to transfer 3 pages; and if by chance
the \kw{repeat} statement covers exactly 5 pages, then there will be a page
fault.

\section{Priorities in the resources used by a program}

The page-management policy, in our opinion, consider that the pages used for
the pseudo-stack and the stack are never preempted, for several reasons:
\begin{enumerate}
    \item The cost of a page removal is higher for a data page than for a code
page because removal of the former implies a copy in the back store; the latter
allows immediate overlay;
    \item If we delete some data pages of a running process, we will have
data-originated fault pages.
\end{enumerate}
On the other hand, we may use the number of data pages and of code pages
requested by a program to establish dynamic priorities between the programs:
for instance, the program time slice may be proportional to its number of code
pages, i.e., we give more time to the programs involved in loop computations;
or we may sort the program waiting queue according to the number of data pages,
i.e., we give priority to the ``memory monopolizing'' programs.

\section{Page size}

In order to respect the philosophy of the system, the page size will be chosen
according to the performance of the page transfer device and of the monitor
program: the transfer time for a page must be less than the average execution
time of the instructions contained in one page (in order to avoid disruption of
a process computation).

This requirement may limit the application of this kind of operating system to
computers using ``slave memory'' (like the {\sc ibm} 360/85 cache store) or
``distributive memory'' (for instance, one level of standard core memory, the
next level being large capacity store LCS, the last one a drum). The previous
discussion explains the kind of policy that would be used to move the pages
between the different levels of the memory hierarchy.

\section{Deadlock}

The approach presented in previous sections must be studied to understand how
it may lead to a deadlock; this deadlock is owing to requests for data pages
when there is no free page frame. But it seems that in this case, preemption of
a job would solve the problem. (Deadlocks due to peripherals are another
problem.)

\section{Conclusion}

This is only an outlook of what may be done about a holon operating system; the
profit of the use of the structured programming methodology would be enhanced
by the improved performance of computer systems.

\section{Some additional comments about the holon operating system}

Hoare and McKeag reject the possibility of prepaging.\autocite{hoare1972survey}
The first reason is the inability to have a significant knowledge of the future
behavior of the programs. The second one is that the rate of a program
execution in a multiprogramming environment is unknown.

Nevertheless, it seems that prepaging through the monitor control can be
feasible with some kind of memory, the latency time of which is small;
moreover, if it appears that a program generates page faults under prepaging,
the monitor may try to avoid this by some modification of the scheduling policy
(for instance, if we have a round-robin policy, we may shorten the time-slice
of the faulting program.)

In the holon system, the synthesizer can use the description of the computer
architecture to choose the place and the extent of the page requests. The
avoidance of page faults is a design goal, not a guarantee.

The Atlas system,\autocite{baylis1968paging} uses an elaborated algorithm to
detect pages involved in loop computations. The Burroughs {\sc b}6700 tries to
define a ``working set'' for each program and adjusts it in order to provide a
fixed rate of page faulting.\autocite{fenton1972b6700} The prepaging policy
that we have described should be compared to these kinds of designs.

%************************************************
\chapter{Conclusion}
%************************************************

In this survey we have described the main features of the holon concept and its
use in programming languages.

As ``the proof of the pudding is in the eating,'' we presently start an
implementation of the holon programming language. With it, it will be possible
to evaluate practically the usefulness of the language.

\bigskip
\bigskip
\bigskip

\hfill Pierre-Arnoul de Marneffe

\hfill Universit\'e de Li\`ege

\hfill Li\`ege Belgium

%%%%%%%%%%%%%%%%%%%%%%%%%%%%%%%%%%%%%%%%%%%%%%%%%%%
%%%%%%%%%%%%%%%%%%%%%%%%%%%%%%%%%%%%%%%%%%%%%%%%%%%
%%%%%%%%%%%%%%%%%%%%%%%%%%%%%%%%%%%%%%%%%%%%%%%%%%%
%%%%%%%%%%%%%%%%%%%%%%%%%%%%%%%%%%%%%%%%%%%%%%%%%%%

\chapter{Appendix}

In this appendix, we comment on a few peculiar points.

\section{The use of the exit statement}

The \kw{exit} statement is only a \kw{goto} the end of a holon; this
restriction forbids the \kw{exit} statement to be used to get out of a loop.
Using the \kw{exit} statement, you first reach the end of the repetitive part
of the loop; but next, you must pass through the test part of the loop.

Example: The searching of a file for the appearance of a given key. There are
two exits of the searching process: either an item matching the key is found,
or the end of the file is reached without finding any matching item.

A holon program for this may be written as follows:

\smallskip\noindent
search for a matching \%KEY\%\\
\ind{1}\kw{begin} \kw{do} try to match next item in the file to \%KEY\%\\
\ind{2}\kw{while} the end of file is not reached or the item is not found;\\
\ind{1}\kw{end}

\smallskip\noindent
the end of file \kw{etc} \\
\ind{1}\kw{begin} \\
\ind{2}\kw{declare} unfound, NOEOF: \kw{boolean at} search \kw{etc}
\kw{level}\\
\ind{2}\kw{initial} \# unfound $\gets$ {\it true}; NOEOF $\gets$ {\it true}
\#\#;\\
\ind{2}(unfound $\equiv$ {\it true}) \kw{and} (NOEOF $\equiv$ {\it true})
\#\#\\
\ind{1}\kw{end}

\smallskip\noindent
try to match \kw{etc} \\
\ind{1}\kw{begin} \\
\ind{2}\kw{declare} F: \kw{file at} search \kw{etc} \kw{level}\\
\ind{2}\kw{declare} index: \kw{integer at} search \kw{etc} \kw{level}\\
\ind{2}\kw{initial} \# index $\gets$ $1$ \#\#;\\
\ind{2}\kw{if} \# F(index) $\equiv$ 'EOF' \#\# \kw{then} \# NOEOF $\gets$ {\it
false}
\#\# \kw{exit};\\
\ind{2}\kw{if} \# F(index) $\equiv$ \%KEY\% \#\#\\
\ind{3}\kw{then} \# unfound $\gets$ {\sl false} \#\# \kw{exit}\\
\ind{3}\kw{else} \# index $\gets$ index $+$ $1$ \#\#;\\
\ind{1}\kw{end}

\smallskip\noindent
\%KEY\% is not declared because it must appear in an outer holon requesting
this search.

As a matter of fact, the example is too simple to show the use of \kw{exit}.
We can code the two \kw{if}s like a \kw{case} statement because there is
nothing to do after the two tests. An obvious way to code the program is
something like:

\medskip
\noindent
\ind{2}index $\gets$ $1$;\\
\ind{2}\kw{while} (F(index) $\neq$ \%KEY\%) \kw{and} NOEOF\\
\ind{2}\kw{do} index $\gets$ index $+$ $1$;

\medskip
\noindent
But the matter here was to show how people must use the \kw{exit} statement,
if one wants to use it.

\section{Understanding the holon concept}

The author really doesn't know to what extent the reader can grasp the holon
concept without reading Koestler's book.\autocite{koestler1968ghost}

Nevertheless, we shall give two other examples of using the holon model. The
first one, quoted from ``The Revolutionist's Handbook'' by George Bernard Shaw,
is a case for the necessity of environment-checking. The second one is a
personal interpretation of some early American building designs; it will
expound the tendency to integration.

\subsection{The self-controlled glutton}

This description was used by Bernard Shaw to expound his statement that: ``The
survival of the fittest means finally the survival of the self-controlled,
because they alone can adapt themselves to the perpetual shifting of
conditions\,\dots .''

``The Glutton, as the man with the strongest motive for nourishing himself,
will always take more pains than his fellows to get food. When food is so
difficult to get that only great exertions can secure a sufficient supply of
it, the Glutton's appetite develops his cunning and enterprise to the utmost;
and he becomes not only the best fed but the ablest man in the community. But
in more hospitable climates, or where the social organization of the food
supply makes it easy for a man to overeat, then the Glutton eats himself out of
health and finally out of existence.''\autocite{shaw1903revolutionist}

\subsection{The integration design at work}

An example of two ways (an unintegrated one and an integrated one) of
implementing the same function can be seen in two early American buildings:
George Washington's Mount Vernon and Thomas Jefferson's Monticello.

It seems that a plantation of that time was composed by a lot of ``units'': the
mansion, the laundry, the smoke house, the kitchen, the dairy, the stable, the
weaving house, etc\,\dots .

Mount Vernon looks like a small town, each of the ``units'' implemented like a
house.

At Monticello, Jefferson put together all the units by locating them beneath
two long terraces forming aisles of the mansion. All the ``units'' are
connected by an all-weather passageway. This way to ``integrate'' the buildings
renders all the ``units,'' except the mansion, as inconspicuous as possible.

The final results, as a visitor can see, are completely different;
nevertheless, they both have fulfilled the same function.

\section{Chronology}

The first report on holon programming circulated privately in December
'72.\autocite{marneffe1972}

In February '73, a submitted draft was accepted by the selection committee of
the International Computing Symposium '73. The ICS was held in Davos,
Switzerland, the 4th through the 7th of September '73.

The paper ``Holon Programming'' will be published in the ICS '73 proceedings by
North-Holland.\autocite{marneffe1973}

In September '73, P. Henderson (University of Newcastle upon Tyne)  made us
aware of a series of three articles on the CADES system by David Pearson (ICL)
published in Computer Weekly (7/26/73, 8/2/73, 8/19/73). The CADES system is
using a top-down approach called ``structural modelling.'' The designers of
this system call a ``holon'' a unit of software design. The CADES system is a
data-oriented one: ``\dots The emphasis is on code being designed to manipulate
defined data structures rather than data structures being invented to support
the codification of abstract functions\,\dots'' That is working in the opposite
direction of our approach. Calling their units ``holons'' is correct because
these units have the two features of a holon: (1) ``\dots A highly resilient,
modular design results with each holon only manipulating a strictly-defined
small subset of the universal set of available data''; and (2) ``\dots It is to
be emphasized that one level only will appear finally in machine-executable
codified form\,\dots''

I am very happy that other people are tempted to use the ``holon concept'' to
master program complexity.

\chapter{Acknowledgments}

First of all, I want to thank Prof. D. Ribbens (University of Li\`ege,
Belgium), who is my teacher in computer science, especially  for the many
discussions and comments about this work.

All our gratitude for Prof. E.W. Dijkstra for his many helpful comments he made
on previous reports.

I want to thank Prof. H. Leroy (Institut d'Informatique, Namur, Belgium), Prof.
E. Milgrom (University of Louvain-la-Neuve, Belgium), and Mr. W.H. Feijen
(T.H.E. Eindhoven) for the useful comments and criticisms they made.

I thank the Belgian National Fund for Scientific Research (F.N.R.S.) and
IBM--Belgium to have granted our stay at the IBM Federal System Center
(Gaithersburg, Maryland).

I am indebted to Dr. Harlan Mills of IBM Gaithersburg, who so kindly accepted
my stay at IBM Gaithersburg. I shall never forget how the people there and
especially Mrs. Carole Hall, Mr. Jay Arnold, Jayadev Misra, Jim Whitlock, and
Bill Patterson welcomed me so friendly.

I must thank Mrs. Barbara Heflin, who managed to type this text so efficiently,
and the Word Processing and Art Department for the work they have done.

\chapter{Transcriber's Endnote}

The phrasings used by Prof. de Marneffe have been preserved except in a few
cases where a gallicism obscured the meaning. Please send any bugs or
suggestions to {\tt holon.programming@gmail.com}.

Thanks to Don Knuth for directing me to the only extant hard copy of this
survey,\marginnote{And, of course, for {\TeX}\ and \mbox{literate
programming}!} for sending a copy of his 1974 letter to the author, and for
allowing the letter to be reproduced here. Thanks to the librarians at the
German National Library of Science and Technology in Hanover, Germany and those
at Love Library in Lincoln, Nebraska who helped me get a hold of this rare
book. Thanks to the creators of the {\tt tufte-book} document class for
creating such a lovely template.

And many thanks to the late Prof. Pierre-Arnoul Frédéric Guy Donat de Marneffe
(1946--2023) for his work and his kind answers to my inquiries.

\bigskip

\hfill Mitch Gerrard

\hfill April 2024

%%
% The back matter contains appendices, bibliographies, indices, glossaries,
% etc.

\backmatter

\chapter{Bibliography}
\printbibliography[heading=none]

\end{document}